\documentclass[11pt,a4paper]{article}
\PassOptionsToPackage{obeyspaces,spaces}{url}

% --- Mode de compilation ---
% Par defaut, compilation rapide (sans microtype, environ 20 % plus rapide).
% Pour la version envoyee au jury ou imprimee, remplacer \versionfinalefalse par \versionfinaletrue.
\newif\ifversionfinale
\versionfinalefalse

% --- Encodage, langue et typographie ---
% (UTF-8 est l'encodage par defaut de LaTeX : inputenc n'est plus necessaire.)
\usepackage[T1]{fontenc}
\usepackage[french]{babel}
\usepackage{lmodern}
\ifversionfinale
  \usepackage{microtype}
\fi

% --- Mise en page et couleurs ---
\usepackage[margin=2.0cm, headheight=15pt]{geometry}
\usepackage[table]{xcolor}
\usepackage{fancyhdr}
\usepackage{titlesec}
\usepackage{parskip}

% --- Tableaux, maths et listes ---
\usepackage{amsmath}
\usepackage{booktabs}
\usepackage{tabularx}
\usepackage{longtable}
\usepackage{enumitem}
\usepackage{float}

% --- Encadres, schemas et code ---
\usepackage[breakable]{tcolorbox}
\usepackage{listings}
\usepackage{tikz}
\usetikzlibrary{arrows.meta, shapes.geometric, fit, backgrounds, babel}

% --- Annexes et liens ---
\usepackage[toc,page]{appendix}
\usepackage[hidelinks]{hyperref}

% --- Palette officielle Crash-Test (version lecture claire / impression) ---
\definecolor{ctPrimary}{HTML}{1B3A5C}
\definecolor{ctMethod}{HTML}{1D6FD8}
\definecolor{ctGenerator}{HTML}{C95E12}
\definecolor{ctCorrection}{HTML}{1F883D}
\definecolor{ctAltOne}{HTML}{8250DF}
\definecolor{ctAltTwo}{HTML}{1B7C83}
\definecolor{ctBaseline}{HTML}{57606A}
\definecolor{ctBroken}{HTML}{CF222E}
\definecolor{ctBgLight}{HTML}{F6F8FA}
\definecolor{ctRowShade}{HTML}{EEF3F8}
\definecolor{ctBorder}{HTML}{D0D7DE}
\definecolor{ctCodeBg}{HTML}{F4F6F8}

% --- En-tetes et pieds de page ---
\pagestyle{fancy}
\fancyhf{}
\fancyhead[L]{\small\color{ctBaseline}\textbf{Crash-Test} : cahier de cadrage}
\fancyhead[R]{\small\color{ctBaseline}Hackathon ENSAI, 30 h}
\fancyfoot[C]{\small\color{ctBaseline}\thepage}
\renewcommand{\headrulewidth}{0.4pt}

% --- Style des sections ---
\titleformat{\section}
  {\Large\bfseries\color{ctPrimary}}
  {\thesection.}{0.5em}{}[\vspace{2pt}\titlerule]
\titleformat{\subsection}
  {\large\bfseries\color{ctPrimary}}
  {\thesubsection}{0.5em}{}
\titleformat{\subsubsection}
  {\normalsize\bfseries\color{ctBaseline}}
  {\thesubsubsection}{0.5em}{}

% --- Encadres reutilisables (un style commun, quatre couleurs) ---
\tcbset{ctcommon/.style={breakable, coltitle=white, fonttitle=\bfseries, boxrule=0.7pt, arc=3pt, left=8pt, right=8pt, top=6pt, bottom=6pt}}
\newtcolorbox{enbref}{ctcommon, colback=ctMethod!6, colframe=ctMethod, title=En bref}
\newtcolorbox{keybox}[1]{ctcommon, colback=ctBgLight, colframe=ctPrimary, title=#1}
\newtcolorbox{rulebox}[1]{ctcommon, colback=ctBroken!5, colframe=ctBroken, title=#1}
\newtcolorbox{okbox}[1]{ctcommon, colback=ctCorrection!6, colframe=ctCorrection, title=#1}

% --- Etiquettes de priorite ---
\newcommand{\pzero}{\colorbox{ctCorrection!18}{\footnotesize\textbf{P0}}}
\newcommand{\pone}{\colorbox{ctMethod!18}{\footnotesize\textbf{P1}}}
\newcommand{\ptwo}{\colorbox{ctBaseline!18}{\footnotesize\textbf{P2}}}

% --- Configuration listings (UTF-8 francais pour pdflatex) ---
\lstdefinelanguage{Python}{morekeywords={def,return,import,from},sensitive=true,morecomment=[l]{\#},morestring=[b]"}
\lstset{
  basicstyle=\ttfamily\footnotesize,
  backgroundcolor=\color{ctCodeBg},
  frame=single,
  rulecolor=\color{ctBorder},
  breaklines=true,
  breakatwhitespace=false,
  tabsize=2,
  showstringspaces=false,
  xleftmargin=4pt,
  xrightmargin=4pt,
  aboveskip=6pt,
  belowskip=6pt,
  commentstyle=\color{ctBaseline}\itshape,
  morecomment=[l]{\#},
  literate=
    {é}{{\'e}}1 {è}{{\`e}}1 {ê}{{\^e}}1 {ë}{{\"e}}1
    {à}{{\`a}}1 {â}{{\^a}}1
    {î}{{\^i}}1 {ï}{{\"i}}1
    {ô}{{\^o}}1
    {ù}{{\`u}}1 {û}{{\^u}}1
    {ç}{{\c{c}}}1 {œ}{{\oe}}1
    {É}{{\'E}}1 {È}{{\`E}}1 {À}{{\`A}}1
    {«}{{\guillemotleft}}1 {»}{{\guillemotright}}1 {’}{{'}}1
}

\newcolumntype{Y}{>{\raggedright\arraybackslash}X}
\newcolumntype{L}[1]{>{\raggedright\arraybackslash}p{#1}}

% --- Fiche d'une mecanique visuelle (annexe F) : une rubrique par ligne ---
\newenvironment{fiche}{\par\small\noindent\tabularx{\textwidth}{@{}L{3.4cm} Y@{}}\toprule}{\bottomrule\endtabularx\par\medskip}

% --- Code en ligne secable (coupures autorisees apres / . _ etc.) ---
% Usage : \code{backend/api/routes_run.py} (underscores ecrits sans antislash)
\DeclareUrlCommand\code{\urlstyle{tt}}

% --- Filet de securite contre les debordements de lignes ---
\setlength{\emergencystretch}{2em}

% --- Styles TikZ communs ---
\tikzset{
  ctbox/.style={rectangle, rounded corners=4pt, draw=ctPrimary, thick, fill=ctBgLight, align=center, font=\small, inner sep=5pt},
  ctstep/.style={rectangle, rounded corners=4pt, draw=ctMethod, thick, fill=white, align=center, font=\small, inner sep=5pt},
  ctgreen/.style={rectangle, rounded corners=4pt, draw=ctCorrection, thick, fill=ctCorrection!8, align=center, font=\small, inner sep=5pt},
  ctarr/.style={-{Stealth[length=2.4mm]}, thick, color=ctPrimary},
  ctnote/.style={font=\footnotesize\itshape, text=ctBaseline, align=center}
}

\begin{document}

% ============================================================================
% PAGE DE GARDE
% ============================================================================
\begin{titlepage}
  \centering
  \vspace*{1.5cm}
  {\Large\color{ctMethod}\textbf{CAHIER DE CADRAGE}}\\[0.5cm]
  {\Huge\bfseries\color{ctPrimary} CRASH-TEST}\\[0.4cm]
  {\large\color{ctBaseline} Casser les hypothèses d'un théorème, mesurer ce qui cède, puis réparer :\\ un laboratoire de simulations et un entraînement par QCM pour les filières quantitatives}\\[1.2cm]

  \begin{tcolorbox}[
    colback=ctBgLight,
    colframe=ctPrimary,
    width=0.94\textwidth,
    arc=4pt,
    boxrule=0.8pt,
    left=12pt, right=12pt, top=10pt, bottom=10pt
  ]
    \textbf{À quoi sert ce document :} c'est la référence commune de l'équipe pendant les 30 heures. Le corps du document explique simplement \emph{ce que l'on construit et pourquoi}, \emph{comment les pièces s'emboîtent} et \emph{qui fait quoi, quand}. Les détails techniques (format YAML complet, contrats d'API, formats des vues, arborescence) sont regroupés en annexes, à consulter au moment de coder. Toutes les décisions sont prises : elles ne se rediscutent pas pendant le hackathon.
  \end{tcolorbox}

  \vspace{1.0cm}
  \begin{tabularx}{0.92\textwidth}{@{}l Y@{}}
    \toprule
    \textbf{Public cible} & Étudiants et enseignants des filières quantitatives (statistiques, probabilités, mathématiques, économie, finance, physique, SII, machine learning). \\
    \midrule
    \textbf{Technologies} & Calcul en Python (\texttt{FastAPI}, NumPy, SciPy, statsmodels), interface en \texttt{Svelte 5} avec Plotly.js et KaTeX, \texttt{Supabase} uniquement pour les liens de partage. \\
    \midrule
    \textbf{Mise en ligne} & Une web app accessible par un simple lien : interface sur \texttt{Cloudflare}, moteur Python sur un hébergeur Python gratuit, version locale en secours. \\
    \midrule
    \textbf{Principe clé} & Aucun code spécifique à un théorème, aucun appel à une IA pendant l'utilisation : chaque expérience est un fichier YAML validé, que l'on peut créer et partager. \\
    \midrule
    \textbf{Équipe} & 4 membres, un binôme de jour et un binôme de nuit en relais (toujours au moins 2 personnes actives, les 4 réunies au démarrage, aux relèves et pour le sprint final). \\
    \bottomrule
  \end{tabularx}

  \vfill
  {\small\color{ctBaseline} Document de référence interne \textbullet{} Version 3.0}
\end{titlepage}

% ============================================================================
% LE PROJET EN UNE PAGE
% ============================================================================
\section*{Le projet en une page}
\addcontentsline{toc}{section}{Le projet en une page}

\begin{keybox}{Le problème}
En cours, un théorème arrive avec une liste d'hypothèses à retenir (variance constante, normalité, indépendance, petits angles, convexité\dots). Les étudiants les récitent sans savoir \emph{ce qui casse} quand l'une d'elles n'est plus vraie, \emph{à partir de quand} la méthode devient trompeuse, ni \emph{quelle correction} la répare. Et quand ils révisent avec des QCM, la correction tient en une ligne de texte : on croit avoir compris, sans l'avoir vu.
\end{keybox}

\begin{okbox}{Notre réponse}
\textbf{Crash-Test, ce sont des QCM de révision où la correction est une expérience.} Pour chaque théorème, l'étudiant peut :
\begin{center}
  \textbf{\color{ctBroken}Casser une hypothèse} $\longrightarrow$ \textbf{\color{ctMethod}Mesurer ce qui cède} $\longrightarrow$ \textbf{\color{ctCorrection}Appliquer la correction}
\end{center}
Deux modes utilisent exactement les mêmes fichiers : le \textbf{Laboratoire} pour comprendre et manipuler, l'\textbf{Entraînement} pour réviser en QCM et mesurer sa lucidité (score de Brier). Un bouton relie les deux.
\end{okbox}

\begin{keybox}{Les décisions qui structurent le projet}
\begin{itemize}[leftmargin=*, nosep]
  \item \textbf{Une seule mécanique pour tous les théorèmes :} une expérience est un fichier YAML, lu par une seule fonction \texttt{run()}. Ajouter un théorème, c'est écrire un fichier, pas du code.
  \item \textbf{Une seule source de vérité :} chaque notion (ce qui est \og rompu \fg{}, l'état sain, la couleur d'un terme) est déclarée une fois dans le YAML ; le moteur en déduit tout le reste.
  \item \textbf{Tout est vérifié au chargement :} une expérience mal écrite échoue avec un numéro de ligne, jamais pendant la démo.
  \item \textbf{Calcul déterministe :} même état et même graine aléatoire donnent toujours le même résultat. C'est fiable, gratuit et reproductible devant le jury.
  \item \textbf{Chaque théorème est expliqué :} énoncé, rôle de chaque hypothèse et conséquence de sa rupture, en deux ou trois lignes.
  \item \textbf{Partage par simple lien,} sans comptes ni galerie communautaire.
  \item \textbf{Périmètre tenu :} 3 expériences irréprochables plutôt que 6 fragiles, un socle P0 resserré pour 30 heures, et une liste de coupes décidée à l'avance.
\end{itemize}
\end{keybox}

\begin{rulebox}{Ce que le jury doit retenir}
En moins de 15 secondes, il voit 100 intervalles de confiance tomber, la couverture s'effondrer autour de 75\,\%, puis revenir près de 95\,\% d'un clic sur Correction. Ensuite, il comprend que la même mécanique couvre la régression, les tests et la physique.
\end{rulebox}

\newpage
\tableofcontents
\newpage

% ============================================================================
% SECTION 1 : LE PROJET EN CLAIR
% ============================================================================
\section{Le projet en clair}

\begin{enbref}
Crash-Test transforme les hypothèses d'un théorème en curseurs. On voit la méthode du cours décrocher, on mesure de combien, puis on la répare. Les étudiants y viennent pour réviser en QCM ; la simulation sert de correction.
\end{enbref}

\subsection{Le constat}
Dans l'enseignement supérieur scientifique, un théorème est presque toujours présenté sous la forme \og résultat + liste d'hypothèses \fg. Trois questions restent pourtant sans réponse concrète pour l'étudiant :
\begin{enumerate}[leftmargin=*, nosep]
  \item \textbf{Qu'est-ce qui cède} quand une hypothèse est relâchée ? (le biais ? la variance ? le niveau du test ?)
  \item \textbf{À partir de quel seuil} la méthode devient-elle trompeuse ?
  \item \textbf{Quelle correction} rétablit la fiabilité, et à quel prix ?
\end{enumerate}

\subsection{L'idée : une boucle unique, pour tous les théorèmes}
Chaque expérience suit la même boucle en trois temps, entièrement déterministe et reproductible grâce à une graine aléatoire (\texttt{seed}) :
\begin{itemize}[leftmargin=*, nosep]
  \item \textbf{\color{ctBroken}Casser :} un curseur, un interrupteur ou un sélecteur de loi relâche une hypothèse.
  \item \textbf{\color{ctMethod}Mesurer :} des indicateurs calculés en direct (taux de couverture, risque réel, erreur de période\dots) montrent l'écart avec la cible, avec un verdict \emph{Fiable}, \emph{Fragile} ou \emph{Rompu}.
  \item \textbf{\color{ctCorrection}Réparer :} le bouton Correction superpose la méthode corrigée du cours (écarts-types HC3, test de Wilcoxon, RK4, formule de Borda\dots) et l'écart se referme. En \pone{}, jusqu'à deux méthodes alternatives peuvent être comparées en plus.
\end{itemize}
Chaque expérience déclare ses méthodes en liste : la méthode du cours, sa correction et d'éventuelles alternatives. Le moteur les calcule toutes sur les mêmes données, à chaque exécution. Afficher ou masquer une méthode ne demande donc aucun nouveau calcul.

\subsection{Pourquoi un étudiant reviendrait sur l'app}
Pour \textbf{réviser avant un partiel}. Le mode Entraînement propose des QCM renouvelés à chaque tirage ; après chaque réponse, la simulation montre pourquoi c'est juste ou faux. C'est la porte d'entrée naturelle, et le Laboratoire est là pour approfondir.

\subsection{Ce qui nous distingue}
\begin{table}[H]
  \centering
  \small
  \begin{tabularx}{\textwidth}{@{}L{4.2cm} Y Y@{}}
    \toprule
    \textbf{Outil existant} & \textbf{Ce qu'il fait bien} & \textbf{Ce qui lui manque} \\
    \midrule
    Banques de QCM, flashcards & S'entraîner vite & Une correction en texte, qui ne montre rien \\
    Seeing Theory, StatPlayground & Visualiser un concept & Pas de QCM, pas de réparation, une seule discipline \\
    ChatGPT et assistants IA & Expliquer à la demande & Rien de vérifié par le calcul, rien à manipuler \\
    \rowcolor{ctRowShade}
    \textbf{Crash-Test} & QCM, simulation et correction réunis, sur plusieurs disciplines & Il mesure en plus la lucidité de l'étudiant (score de Brier) \\
    \bottomrule
  \end{tabularx}
  \caption{Positionnement de Crash-Test face aux outils existants.}
\end{table}

% ============================================================================
% SECTION 2 : CE QUE VIT L'ETUDIANT
% ============================================================================
\section{Ce que vit l'étudiant}

\begin{enbref}
Deux modes, les mêmes fichiers. Le Laboratoire sert à comprendre et à manipuler librement ; l'Entraînement sert à réviser en QCM. Après chaque réponse, un bouton ouvre le Laboratoire exactement sur la situation de la question.
\end{enbref}

\begin{figure}[H]
  \centering
  \begin{tikzpicture}[
    col/.style={ctstep, text width=4.9cm, minimum height=1.0cm},
    tcol/.style={ctstep, draw=ctGenerator, text width=4.9cm, minimum height=1.0cm},
  ]
    % Accueil
    \node[ctbox, text width=4.2cm, minimum height=0.9cm] (home) at (0,0) {\textbf{Accueil}\\ ou lien de partage reçu};

    % Colonne Laboratoire
    \node[font=\bfseries\color{ctMethod}] (labT) at (-4.2,-1.25) {Laboratoire : comprendre};
    \node[col] (l1) at (-4.2,-2.1) {Choisir une expérience\\dans la galerie};
    \node[col] (l2) at (-4.2,-3.6) {\textbf{Comprendre :} l'énoncé et\\ses hypothèses cliquables};
    \node[col] (l3) at (-4.2,-5.1) {\textbf{Casser :} curseurs,\\interrupteurs, lois};
    \node[col, draw=ctCorrection] (l4) at (-4.2,-6.6) {\textbf{Mesurer et réparer :}\\la Correction referme l'écart};
    \node[col] (l5) at (-4.2,-8.1) {Défis, puis partage\\par lien};

    % Colonne Entrainement
    \node[font=\bfseries\color{ctGenerator}] (trT) at (4.2,-1.25) {Entraînement : réviser};
    \node[tcol] (t1) at (4.2,-2.1) {Question générée :\\graphique ou texte};
    \node[tcol] (t2) at (4.2,-3.6) {\textbf{Confiance :} Je devine,\\Plutôt sûr ou Certain};
    \node[tcol] (t3) at (4.2,-5.1) {\textbf{Validation par simulation}\\+ explication};
    \node[tcol] (t4) at (4.2,-6.6) {\textbf{Bilan :} score de Brier\\et courbe de calibration};

    % Fleches
    \draw[ctarr] (home.west) -| (labT.north);
    \draw[ctarr] (home.east) -| (trT.north);
    \draw[ctarr] (l1) -- (l2);
    \draw[ctarr] (l2) -- (l3);
    \draw[ctarr] (l3) -- (l4);
    \draw[ctarr] (l4) -- (l5);
    \draw[ctarr] (t1) -- (t2);
    \draw[ctarr] (t2) -- (t3);
    \draw[ctarr] (t3) -- node[right, ctnote] {après la série} (t4);
    \draw[ctarr, ctMethod, very thick] (t3.west) -- node[above, font=\footnotesize\bfseries, text=ctMethod, align=center] {Ouvrir dans\\le Laboratoire} node[below, ctnote] {sur l'état exact} (l3.east);
    \draw[ctarr, dashed] (t3.north west) .. controls +(-1.2,0.9) and +(1.6,0) .. node[above, ctnote, pos=0.5] {revoir l'hypothèse} (l2.east);
  \end{tikzpicture}
  \caption{Le parcours de l'étudiant : deux portes d'entrée reliées par un pont.}
\end{figure}

\subsection{Le Laboratoire : comprendre et manipuler librement}
\begin{itemize}[leftmargin=*]
  \item \textbf{Accès immédiat :} l'étudiant choisit une expérience dans la galerie (filtrable par discipline et par famille d'hypothèses) et accède tout de suite aux contrôles. Aucune question ne bloque l'écran.
  \item \textbf{Le panneau \og Comprendre \fg{} :} il affiche l'énoncé du théorème et la liste de ses hypothèses. Un clic sur une hypothèse la casse dans la simulation et affiche deux ou trois lignes : \emph{à quoi elle sert} et \emph{ce que l'on voit quand elle tombe}. C'est ce qui fait de Crash-Test un outil qui explique, et pas seulement qui montre (détails en section~\ref{sec:explications}).
  \item \textbf{Une seule grande vue,} organisée en onglets (Simulation, Résidus, Répétitions, Étude de sensibilité, Domaine de validité), pour rester lisible.
  \item \textbf{Des contrôles adaptés à chaque hypothèse :} curseur continu $h \in [0,1]$ avec raccourcis 0 et 1, interrupteur \emph{Respectée / Relâchée}, ou sélecteur de loi (Normale, Uniforme, Exponentielle, Student, Cauchy, Bimodale\dots). On peut casser plusieurs hypothèses à la fois et charger des contre-exemples en un clic. Un contrôle qui n'a pas de sens dans l'état courant est grisé avec une phrase d'explication (champ \code{disabled_when}) : aucun état absurde n'est calculé.
  \item \textbf{Un verdict lisible sur chaque mesure :} \emph{Fiable}, \emph{Fragile} ou \emph{Rompu}, avec la marge due au hasard des tirages (\og 74,0\,\% $\pm$ 2,0 \fg{}).
  \item \textbf{Le bouton Correction} \pzero{} superpose la méthode corrigée du cours et montre l'écart qui se referme.
  \item \textbf{Le banc de méthodes} \pone{} : jusqu'à deux méthodes alternatives s'affichent à côté de la correction (par exemple MCO, HC3 et Newey-West, ou moyenne et médiane). Activé à $h = 0$, le banc montre le prix d'une méthode robuste quand tout va bien : intervalles plus larges, puissance plus faible, ou biais.
  \item \textbf{L'onglet Défis} propose 2 ou 3 micro-objectifs vérifiés par le calcul, pour donner une raison de toucher aux curseurs (exemple : \og piégez le test de Breusch-Pagan \fg{}).
  \item \textbf{Une question d'ouverture facultative} \pone{} : une carte \og Pariez avant de toucher aux curseurs \fg{} avec un bouton \emph{Passer}. Elle garde l'esprit \og prédire, observer, expliquer \fg{} sans bloquer personne.
\end{itemize}

\subsection{L'Entraînement : réviser et mesurer sa lucidité}
\begin{itemize}[leftmargin=*]
  \item \textbf{Séries de 10 à 15 questions,} ou mode continu, filtrables par discipline ou par expérience.
  \item \textbf{Des questions toujours nouvelles :} chaque tirage change la graine, les paramètres numériques et l'ordre des choix. Impossible d'apprendre une réponse ou un graphique par cœur.
  \item \textbf{Trois formats :} les idées fausses classiques (\og l'hétéroscédasticité biaise la pente ? \fg{}), les diagnostics sur graphique générés à la volée (\og quelle pathologie voyez-vous sur ces résidus ? \fg{}) et, en \pone{}, les prédictions (\og que devient la couverture si $n$ double ? \fg{}), dont la bonne réponse est calculée par le moteur et jamais écrite à la main.
  \item \textbf{Des séries prêtes et partageables :} une série s'ouvre immédiatement, car le serveur en garde quelques-unes d'avance et fabrique les suivantes en arrière-plan. Chaque série a son lien : deux étudiants qui ouvrent le même lien ont les mêmes questions dans le même ordre, et peuvent comparer leurs Bilans (section~\ref{sec:fabrique}).
  \item \textbf{Des corrections toujours justes :} chaque état de validation porte une condition (\texttt{verify}) que la simulation doit confirmer. Si les paramètres tirés au hasard la rendent fausse, le serveur tire d'autres paramètres ; la question n'est jamais servie avec une correction contredite par le calcul.
  \item \textbf{Des diagnostics sans saisie fastidieuse :} les pathologies proposées sont fabriquées automatiquement à partir des hypothèses de l'expérience (une option de loi, un interrupteur cassé, un curseur au maximum), avec leurs libellés et leurs explications.
  \item \textbf{La validation par simulation, toujours :} après chaque réponse, juste ou fausse, la simulation tourne sur l'état qui confirme le raisonnement ou contredit l'erreur, accompagnée de l'explication de l'hypothèse concernée.
  \item \textbf{Le pont vers le Laboratoire :} le bouton \og Ouvrir dans le Laboratoire \fg{} déverrouille tous les contrôles sur cet état précis, avec le panneau Comprendre ouvert sur la bonne hypothèse.
  \item \textbf{Trois crans de confiance} et un \textbf{écran Bilan} séparé, présentés ci-dessous.
\end{itemize}

\subsection{La confiance et le score de Brier, expliqués simplement}

\textbf{L'idée.} Avant de valider, l'étudiant indique s'il est sûr de lui. On ne mesure pas seulement s'il a juste, mais \emph{s'il sait quand il a juste}. C'est la mesure chiffrée de l'illusion de compétence.

\textbf{Trois boutons, toujours les mêmes.} À l'écran, l'étudiant voit trois mots. Derrière, la probabilité dépend du nombre de choix $k$ de la question :

\begin{table}[H]
  \centering
  \small
  \begin{tabularx}{0.9\textwidth}{@{}l Y c c@{}}
    \toprule
    \textbf{Bouton} & \textbf{Probabilité associée} & \textbf{Binaire ($k=2$)} & \textbf{4 choix ($k=4$)} \\
    \midrule
    \og Je devine \fg{} & le hasard, $1/k$ & 50\,\% & 25\,\% \\
    \og Plutôt sûr \fg{} & le milieu entre le hasard et 95\,\% & 72,5\,\% & 60\,\% \\
    \og Certain \fg{} & 95\,\% & 95\,\% & 95\,\% \\
    \bottomrule
  \end{tabularx}
  \caption{Les trois crans de confiance et leur traduction en probabilité.}
\end{table}

Pourquoi trois crans et pas quatre : la décision est plus rapide, chaque cran reçoit plus de réponses (la courbe est donc moins bruitée), et on évite un cran à 25\,\% absurde sur une question binaire. Pourquoi 95\,\% et pas 100\,\% pour \og Certain \fg{} : une seule erreur à 100\,\% écraserait tout le score, et on n'est jamais totalement sûr.

\textbf{Le score de Brier.} Pour chaque question, on compare la confiance annoncée $p$ au résultat $y$ (1 si juste, 0 si faux) :
\[
  \mathrm{Brier} = \frac{1}{N} \sum_{i=1}^{N} (p_i - y_i)^2 .
\]
Il va de 0 (parfait) à 1 (le pire). La pénalité est asymétrique : être \og Certain \fg{} et se tromper coûte $0{,}95^2 \approx 0{,}90$, alors que dire \og Je devine \fg{} limite les dégâts.

\begin{keybox}{L'exemple qui parle au jury}
Sur des QCM à 4 choix, deux étudiants ont chacun 60\,\% de bonnes réponses. Le premier clique toujours \og Plutôt sûr \fg{} (60\,\%) : son Brier vaut \textbf{0,24}. Le second clique toujours \og Certain \fg{} (95\,\%) : son Brier vaut \textbf{0,36}. Ils savent autant de choses, mais le second est nettement pénalisé pour son excès de confiance.
\end{keybox}

\textbf{L'écran Bilan.} Pour ne pas charger l'écran de question, l'historique est regroupé dans un onglet à part :
\begin{itemize}[leftmargin=*, nosep]
  \item une phrase lisible avant tout chiffre : \og Quand vous dites \emph{Certain}, vous avez raison dans 60\,\% des cas \fg{} ;
  \item la \textbf{courbe de calibration} (taux de réussite réel par bouton de confiance), affichée à partir d'une vingtaine de réponses ; avant, un simple \og encore $X$ questions \fg{} ;
  \item le score de Brier.
\end{itemize}
Pour la soutenance, un bouton discret charge un \textbf{profil de démonstration} d'une trentaine de réponses, construit pour montrer clairement un étudiant trop sûr de lui. Le détail mathématique (dont la décomposition de Murphy) est en annexe~\ref{ann:brier}.

\subsection{Expliquer les théorèmes, pas seulement les casser}\label{sec:explications}
Chaque expérience contient de courts textes d'explication, écrits à l'avance dans le fichier YAML (aucune IA à l'exécution) :

\begin{table}[H]
  \centering
  \small
  \begin{tabularx}{\textwidth}{@{}L{3.6cm} Y L{4.6cm}@{}}
    \toprule
    \textbf{Champ} & \textbf{Ce qu'il contient} & \textbf{Où il s'affiche} \\
    \midrule
    \texttt{theorem} & Le nom et l'énoncé du théorème, en 2 ou 3 lignes & En tête du panneau Comprendre \\
    \code{statement_tex} & L'hypothèse écrite en formule & Dans la liste des hypothèses \\
    \texttt{why} & À quoi sert l'hypothèse dans la preuve ou dans la formule & Au clic sur l'hypothèse \\
    \code{when_broken} & Ce que l'on observe quand elle tombe & Au clic, et après une réponse QCM \\
    \code{methods[].explain} & Pourquoi la correction ou l'alternative répare & À l'affichage de la méthode \\
    \code{measures.texts} & Une phrase par verdict (Fiable, Fragile, Rompu) & Sous la mesure, en direct \\
    \code{disabled_text} & Pourquoi un contrôle est grisé & À côté du contrôle \\
    \bottomrule
  \end{tabularx}
  \caption{Les champs d'explication d'une expérience.}
\end{table}

Exemple pour l'homoscédasticité en régression : \emph{À quoi elle sert :} la formule de variance de la pente suppose un même $\sigma^2$ pour toutes les observations. \emph{Quand elle tombe :} la pente reste juste en moyenne, mais les intervalles sont trop étroits et couvrent moins de 95\,\%.

Le test automatique des expériences vérifie qu'aucune hypothèse et aucune méthode ne sont sans explication.

\subsection{Le partage par lien}\label{sec:partage}
Le partage se limite à un lien, sans comptes ni galerie communautaire :
\begin{itemize}[leftmargin=*, nosep]
  \item \textbf{Pour une expérience officielle,} l'état complet tient dans une URL lisible, que l'on peut taper en démo :\\ \code{?exp=regression-lineaire&v=1&p.n=300&h.hetero=0.85&seed=42&on=hc3}.\\ Les points ajoutés à la main s'écrivent \code{pts=4.2,-2.0;5.1,3.3}. On peut ainsi envoyer un contre-exemple précis à un camarade.
  \item \textbf{Pour une expérience importée} \pone{}, son contenu est enregistré dans une table Supabase sous un identifiant aléatoire, lisible par toute personne qui a le lien : \code{?share=<id>&...}. Une ligne enregistrée ne change plus jamais. Le mécanisme est repris d'une autre application de l'équipe (environ 15 minutes de travail).
  \item \textbf{Chaque lien porte la version du format} (\texttt{v=1}). Un lien d'une version que le moteur ne connaît pas est refusé avec un message clair, au lieu de produire un résultat faux.
\end{itemize}
L'état se sépare en deux parties. L'\textbf{état de calcul} (paramètres, hypothèses, points ajoutés, graine) est envoyé au moteur et détermine le résultat. L'\textbf{état d'affichage} (méthodes affichées, onglet, calques) ne quitte jamais l'interface. Les deux tiennent dans l'URL.

L'historique d'Entraînement reste dans le navigateur (\texttt{localStorage}) : aucune inscription n'est nécessaire, ni pour l'étudiant ni pour le jury.

\subsection{Ton et identité visuelle}
\begin{rulebox}{Trois règles impératives}
  \begin{itemize}[leftmargin=*, nosep]
    \item \textbf{Vouvoiement partout :} interface, énoncés, défis, verdicts, explications et messages d'erreur (\emph{\og Vous observez\dots \fg}, \emph{\og Modifiez\dots \fg}).
    \item \textbf{Thème sombre scientifique} défini dans \texttt{tokens.css}, avec une palette de sens unique sur les équations, les badges et les graphiques :
    \begin{center}
      \small
      \textcolor{ctMethod}{\textbf{\#4ea1ff} Méthode} \quad|\quad
      \textcolor{ctCorrection}{\textbf{\#3fb950} Correction} \quad|\quad
      \textcolor{ctAltOne}{\textbf{\#d2a8ff} Alternative 1} \quad|\quad
      \textcolor{ctAltTwo}{\textbf{\#56d4dd} Alternative 2} \\[2pt]
      \textcolor{ctGenerator}{\textbf{\#f0883e} Générateur} \quad|\quad
      \textcolor{ctBaseline}{\textbf{\#6e7681} Référence} \quad|\quad
      \textcolor{ctBroken}{\textbf{\#f85149} Rupture}
    \end{center}
    \item \textbf{Jamais la couleur seule :} chaque rôle a aussi un style de trait et un marqueur (pointillés pour la Référence, losanges pour la Correction, croix pour les échecs\dots), pour rester lisible par une personne qui distingue mal le rouge du vert. Le tableau complet est en annexe~\ref{ann:vues}.
  \end{itemize}
\end{rulebox}

% ============================================================================
% SECTION 3 : LA BOITE A OUTILS VISUELLE
% ============================================================================
\section{La boîte à outils visuelle}

\begin{enbref}
Quatre vues communes à toutes les expériences, un verdict à trois états, une équation dont les termes changent de couleur, des défis, et sept mécaniques visuelles. Chaque élément porte une priorité : \pzero{} indispensable, \pone{} si le P0 est vert, \ptwo{} après le hackathon.
\end{enbref}

\subsection{Les 4 vues communes à toutes les expériences}
En plus des vues propres à chaque méthode (nuage de points, résidus, histogramme, trajectoire), le moteur sait produire quatre vues transversales :
\begin{enumerate}[leftmargin=*]
  \item \textbf{Référence} \pzero{} : l'état sain reste visible en gris pointillé derrière l'état cassé. L'état sain est défini sans ambiguïté : toutes les hypothèses à leur valeur saine (\code{healthy_value}), aucun point ajouté, mêmes paramètres et même graine. Il ne dépend donc pas des curseurs d'hypothèses : il est calculé une fois par jeu de paramètres et reste en cache.
  \item \textbf{Répétitions} \pzero{} : 100 intervalles de confiance empilés autour de la vraie valeur (les échecs en rouge), ou un faisceau de 30 trajectoires de $\bar{X}_n$. La définition fréquentiste devient visible. Les mesures sont toujours calculées sur toutes les répétitions (500 par défaut), seul l'affichage est limité.
  \item \textbf{Étude de sensibilité} \pzero{} : la mesure de fiabilité en fonction de la perturbation, avec la bande de tolérance, le \textbf{seuil critique} et la courbe corrigée en vert. Par défaut, 11 valeurs de $h$ (de 0 à 1 par pas de 0,1). Pour une hypothèse à choix (la loi des erreurs) ou un interrupteur, un diagramme en barres compare chaque option à la cible.
  \item \textbf{Domaine de validité} \pone{} : une carte de chaleur qui croise deux grandeurs (par exemple la violation et la taille d'échantillon $n$), sur une grille d'au plus $12 \times 12$ cases. Elle montre si un grand échantillon compense le problème (intervalle de Student pour une moyenne, sous une loi asymétrique) ou ne répare rien (hétéroscédasticité en MCO). Un clic sur une case règle les deux contrôles.
\end{enumerate}
La Sensibilité et le Domaine de validité reposent sur une seule opération du moteur, le \textbf{balayage} (\texttt{sweep}) : relancer toutes les méthodes sur une grille d'états, à 1 axe ou à 2 axes. Le bouton $n \to \infty$, le balayage automatique, la vérification des défis et les questions de prédiction utilisent la même opération (annexe~\ref{ann:api}).

\textbf{Trois règles communes à toutes les vues :}
\begin{itemize}[leftmargin=*, nosep]
  \item \textbf{Des axes qui ne bougent pas :} les bornes sont calculées sur la Référence (quantiles 0,5\,\% et 99,5\,\%, plus 10\,\% de marge) et ne changent que si les paramètres ou la graine changent, jamais quand on bouge une hypothèse. L'effet du curseur reste visible, et une loi de Cauchy n'écrase plus le graphique. Les points hors cadre sont dessinés au bord avec un triangle et comptés (\og 3 points hors cadre \fg{}).
  \item \textbf{La marge du hasard affichée :} chaque mesure est accompagnée de son erreur Monte Carlo ; l'interface affiche $\pm$ 2 erreurs types. Pour une couverture de 95\,\% sur 500 tirages, cela fait environ $\pm$ 2 points.
  \item \textbf{Un affichage borné par le moteur :} 100 intervalles, 30 trajectoires, 2\,000 points par nuage et par courbe au plus. Chaque vue indique ce qui est affiché et ce qui a été calculé.
\end{itemize}

\subsection{Le verdict, l'équation dynamique et les défis}\label{sec:verdict}
\begin{itemize}[leftmargin=*]
  \item \textbf{Le verdict} \pzero{} : chaque mesure déclare sa cible, le sens dans lequel l'écart est grave, et deux écarts : celui jusqu'auquel la méthode reste fiable (\code{ok_within}) et celui au-delà duquel elle est rompue (\code{broken_beyond}). Le moteur en déduit le verdict, \emph{Fiable}, \emph{Fragile} ou \emph{Rompu}. Pour la couverture d'un IC à 95\,\% : Fiable à partir de 92\,\%, Rompu sous 85\,\%, Fragile entre les deux. Quand une mesure n'a pas de valeur cible absolue (la puissance d'un test dépend de $n$ et de l'effet), elle déclare \code{target: base} : la cible est alors la valeur de la même méthode à l'état sain, et le verdict mesure ce que la rupture a coûté. La \textbf{mesure primaire} d'une expérience pilote le badge, la couleur de l'équation, la bande de la sensibilité, le seuil critique et la pause du balayage : ces éléments ne peuvent donc pas se contredire.
  \item \textbf{Équation dynamique} \pzero{} : chaque expérience écrit sa formule clé comme un gabarit dont certains termes sont nommés (par exemple $\hat{\sigma}^2 (X^\top X)^{-1}$ dans la variance des MCO, ou l'erreur $O(\Delta t)$ de la méthode d'Euler, qui devient $O(\Delta t^4)$ avec RK4). Chaque hypothèse déclare les termes qu'elle met en cause (champ \texttt{affects}). Ainsi, quand plusieurs hypothèses sont cassées, seul le terme réellement fautif change de couleur, selon la règle ci-dessous.
  \item \textbf{Défis et angles morts des tests} \pzero{} : 2 ou 3 objectifs par expérience, vérifiés par le calcul. Le plus parlant : sur l'échantillon affiché, obtenir une $p$-value de Breusch-Pagan au-dessus de 0,05 (donc rassurante) alors que la couverture réelle est déjà rompue, sous 85\,\%. Il faut un petit $n$ et quelques tirages : à $n = 30$ et $h = 0{,}7$, environ un échantillon sur cinq trompe le test (vérifié par simulation). Le message : ne pas rejeter $H_0$ ne prouve pas que l'hypothèse tient.
\end{itemize}

\begin{table}[H]
  \centering
  \small
  \begin{tabularx}{\textwidth}{@{}Y L{6.2cm}@{}}
    \toprule
    \textbf{Situation} & \textbf{Affichage du terme} \\
    \midrule
    Aucune hypothèse qui le met en cause n'est cassée & version saine, couleur neutre \\
    Une hypothèse cassée, mais la méthode principale reste Fiable & version saine, neutre (la méthode résiste à cette rupture) \\
    Une hypothèse cassée et la méthode principale Fragile ou Rompue & version saine, \textcolor{ctBroken}{\textbf{en rouge}} \\
    Même cas, Correction affichée et Fiable & version corrigée, \textcolor{ctCorrection}{\textbf{en vert}} \\
    Même cas, Correction affichée mais pas Fiable & version corrigée, \textcolor{ctBroken}{\textbf{en rouge}} \\
    \bottomrule
  \end{tabularx}
  \caption{La règle de couleur d'un terme de l'équation. Le rouge commence au verdict Fragile, donc exactement au seuil critique.}
\end{table}

\subsection{Les 7 mécaniques visuelles, par priorité}\label{sec:mecaniques}
Toutes reposent sur le même contrat et fonctionnent dans plusieurs disciplines sans code dédié.

\textbf{Comment lire cette section.} Le grand tableau plus bas dit \emph{à quoi sert} chaque mécanique. Tout ce qu'il faut pour la \emph{coder} est dans sa fiche, en annexe~\ref{ann:mecaniques}, toujours avec les mêmes cinq rubriques : ce que voit l'étudiant, le champ YAML qui l'active, les données que renvoie le moteur, ce que fait l'interface, et les règles et cas limites. Le tableau ci-dessous indique où chercher.

\begin{table}[H]
  \centering
  \small
  \begin{tabularx}{\textwidth}{@{}L{3.9cm} c c Y L{3.4cm}@{}}
    \toprule
    \textbf{Mécanique} & \textbf{Prio.} & \textbf{Fiche} & \textbf{Activée dans le YAML par} & \textbf{Composant Svelte} \\
    \midrule
    Banc de méthodes & \pone & \ref{fiche:banc} & la liste \texttt{methods} & \texttt{MethodBench} \\
    Contre-exemples en un clic & \pzero & \ref{fiche:contre} & \texttt{counterexamples} & \texttt{CounterexampleBar} \\
    Balayage automatique & \pone & \ref{fiche:balayage} & \code{auto_sweep} & \texttt{AutoSweepPlayer} \\
    Ajout de points par clic & \pone & \ref{fiche:points} & \code{allow_point_injection} (vue) et \texttt{inject} (générateur) & \texttt{PlotlyChart} \\
    Journal des tirages & \pone & \ref{fiche:journal} & \code{draw_log} & \texttt{DrawLog} \\
    Inspecteur d'observation & \ptwo & \ref{fiche:inspecteur} & \texttt{inspect} (vue) & \texttt{Inspector} \\
    Calques pas à pas & \ptwo & \ref{fiche:calques} & \texttt{layers} (vue) et \texttt{layer} dans \texttt{bind} & \texttt{LayerStepper} \\
    \bottomrule
  \end{tabularx}
  \caption{Où trouver chaque mécanique : sa fiche en annexe~\ref{ann:mecaniques}, le champ qui l'active et le composant qui l'affiche.}
\end{table}

\begin{table}[H]
  \centering
  \small
  \begin{tabularx}{\textwidth}{@{}L{3.9cm} c Y@{}}
    \toprule
    \textbf{Mécanique} & \textbf{Prio.} & \textbf{Ce qu'elle apporte} \\
    \midrule
    Banc de méthodes & \pone & La Correction et jusqu'à deux alternatives côte à côte (MCO, HC3 et Newey-West ; moyenne et médiane ; Student et Wilcoxon). Activé à $h=0$, il montre qu'une méthode robuste a un prix : variance de la médiane, puissance de Wilcoxon, biais de Ridge. Aucun code dédié : toutes les méthodes sont déjà calculées. \\
    Contre-exemples en un clic & \pzero & Des préréglages écrits dans le YAML. En \pone{}, l'effet Anscombe, généré par le moteur : le même nuage linéaire, puis une relation courbe, puis un point levier, avec une pente et un $R^2$ calés à l'identique par le moteur, mais des résidus opposés. \\
    Balayage automatique & \pone & Un bouton Lecture anime $h$ de 0 à 1, marque une pause au seuil critique, puis active la Correction. Son jumeau parcourt les valeurs de balayage de $n$ (par exemple 15, 30, 100, 500, 2000) et sépare ce que le grand échantillon répare (couverture d'un intervalle de Student sur loi asymétrique, précision d'un estimateur convergent) de ce qu'il ne répare jamais (couverture des MCO classiques sous hétéroscédasticité, biais de variable omise, moyenne empirique sous Cauchy). \\
    Ajout de points par clic & \pone & Un clic dans le nuage pose un point (un clic sur un point posé le retire) : on distingue un point levier (qui fait pivoter la droite) d'un outlier vertical (qui dégrade le $R^2$ sans bouger la pente). La même mécanique sert partout où un générateur l'autorise, par exemple une valeur extrême posée sur un histogramme pour comparer moyenne et médiane. \\
    Journal des tirages & \pone & Des pastilles comptent les tirages successifs (\og 15 tirages sous $H_0$ : 1 faux positif \fg{}) pour rendre concrets le risque $\alpha$ et le p-hacking. Une pastille est le test sur un seul jeu de données, comme dans la vraie vie. \\
    Inspecteur d'observation & \ptwo & Un clic sur un point affiche sa fiche : prédiction, résidu, levier $h_{ii}$, distance de Cook, part dans la somme des carrés des résidus. \\
    Calques pas à pas & \ptwo & Construire une figure couche par couche, sans nouveau calcul. Régression et ANOVA : nuage, puis écarts à la moyenne ($SCT$), puis écarts expliqués ($SCE$), puis résidus ($SCR$), avec $SCT = SCE + SCR$ vérifié en direct. Tests : loi sous $H_0$, puis zones de rejet, puis $p$-value, puis loi sous $H_1$ avec $\beta$ et la puissance. \\
    \bottomrule
  \end{tabularx}
  \caption{Les mécaniques visuelles et leur priorité.}
\end{table}

% ============================================================================
% SECTION 4 : COMMENT CA MARCHE
% ============================================================================
\section{Comment ça marche}

\begin{enbref}
Une expérience est un fichier YAML. Une seule fonction, \texttt{run()}, le lit avec l'état choisi par l'étudiant (curseurs, graine), génère les données, applique toutes les méthodes déclarées, agrège les mesures et renvoie verdicts et graphiques. Chaque changement de l'état de calcul relance \texttt{run()} ; afficher ou masquer une méthode ne relance rien.
\end{enbref}

\subsection{Le principe : aucun code propre à un théorème}
C'est la règle qui rend le projet extensible : \textbf{il n'y a jamais de \code{if experiment.id == "..."}}, ni dans le moteur Python, ni dans l'interface. Tout ce qui est propre à un théorème est décrit dans son fichier YAML, validé par un schéma strict (Pydantic), puis exécuté par une fonction unique \code{run(experiment, state) -> RunResult}.

Au chargement, le moteur vérifie tout ce qui peut l'être : les noms utilisés dans les expressions, l'absence de cycle entre variables, les formes des tableaux, que chaque vue lit des données qu'une méthode fournit, que chaque terme de l'équation est cité par une hypothèse, et les plafonds de calcul (section~\ref{sec:robustesse}).

\begin{figure}[H]
  \centering
  \begin{tikzpicture}[
    inb/.style={ctstep, text width=4.8cm, minimum height=1.05cm},
    side/.style={ctbox, text width=3.1cm, minimum height=1.25cm},
    outbox/.style={ctgreen, text width=3.3cm, minimum height=1.25cm},
  ]
    % Entrees
    \node[side] (yaml) at (0,1.4) {\textbf{Expérience}\\fichier YAML\\(ce qu'on simule)};
    \node[side] (state) at (0,-1.4) {\textbf{État de calcul}\\curseurs, graine\\(ce que l'étudiant a réglé)};

    % Etapes internes de run()
    \node[font=\small\bfseries, text=ctPrimary] (runT) at (5.8,2.85) {\texttt{run(expérience, état)}};
    \node[inb] (s1) at (5.8,1.7) {\textbf{1. Générer le monde}\\{\footnotesize Générateurs : \texttt{random}, \texttt{function}, \texttt{ode}}};
    \node[inb] (s2) at (5.8,0) {\textbf{2. Appliquer toutes les méthodes}\\{\footnotesize principale, correction, alternatives}};
    \node[inb] (s3) at (5.8,-1.7) {\textbf{3. Agréger et juger}\\{\footnotesize mesures, verdicts, défis, vues}};
    \draw[ctarr] (s1) -- (s2);
    \draw[ctarr] (s2) -- (s3);

    % Cadre de run()
    \begin{scope}[on background layer]
      \node[draw=ctPrimary, very thick, rounded corners=6pt, fill=ctPrimary!6, fit=(runT)(s1)(s2)(s3), inner sep=8pt] (run) {};
    \end{scope}

    % Sorties
    \node[outbox] (json) at (11.4,1.4) {\textbf{Résultat (JSON)}\\mesures, verdicts,\\défis, vues};
    \node[outbox] (ui) at (11.4,-1.4) {\textbf{Interface}\\graphiques Plotly,\\équation KaTeX};

    % Fleches
    \draw[ctarr] (yaml.east) -- (yaml.east -| run.west);
    \draw[ctarr] (state.east) -- (state.east -| run.west);
    \draw[ctarr] (json.west -| run.east) -- (json.west);
    \draw[ctarr] (json) -- (ui);

    % Boucle de l'etudiant
    \draw[ctarr, draw=ctCorrection, color=ctCorrection] (ui.south) -- ++(0,-1.25) -| node[pos=0.25, below, font=\footnotesize\itshape, text=ctCorrection] {chaque changement de curseur ou de graine relance \texttt{run()}} (state.south);
  \end{tikzpicture}
  \caption{Le pipeline de Crash-Test : un fichier, une fonction, et une boucle pilotée par l'étudiant.}
\end{figure}

\subsection{Les quatre briques du moteur}
\begin{enumerate}[leftmargin=*]
  \item \textbf{Les Générateurs} produisent le \og vrai monde \fg{} :
  \begin{itemize}[nosep]
    \item \texttt{random} \pzero{} : Monte Carlo vectorisé (500 échantillons d'un coup). Une loi ou une forme de fonction peut être remplacée par une autre depuis le YAML (syntaxe \texttt{@h.<id>}).
    \item \texttt{function} \ptwo{} (format décrit dès maintenant, code écrit avec l'expérience qui en a besoin) : courbes et surfaces sur une grille, avec trois niveaux de liberté pour $f(x)$ : une fonction par défaut, 3 ou 4 préréglages, ou une saisie libre \pone{} à laquelle l'expérience ajoute la perturbation.
    \item \texttt{ode} \pzero{} : équations différentielles ; le générateur fournit la solution de référence (\code{solve_ivp}) et le système, que les méthodes intègrent avec leur propre schéma (annexe~\ref{ann:yaml}).
    \item \texttt{process} \ptwo{} : récurrences discrètes pour les séries temporelles.
  \end{itemize}
  \item \textbf{Les Méthodes} sont les outils du cours et leurs corrections : des fonctions Python pures (MCO, HC3 et Newey-West, Student et Wilcoxon, moyenne et médiane, Euler et RK4\dots), enregistrées dans un registre par un décorateur \code{@register_method("nom")}. Chaque méthode déclare les résultats qu'elle fournit et son coût de calcul (annexe~\ref{ann:methodes}). Elles sont écrites en NumPy par formules fermées, sur les 500 échantillons d'un coup : une boucle \texttt{statsmodels} sur 500 échantillons coûterait plusieurs centaines de millisecondes par méthode. \texttt{statsmodels} et SciPy servent de référence dans les tests, pour vérifier les formules.
  \item \textbf{Les Agrégateurs} transforment les résultats de chaque répétition en mesures : taux de couverture, biais, risque réel, erreur quadratique\dots{} Une méthode ne calcule jamais elle-même une couverture ou un taux de rejet.
  \item \textbf{Les Vues} sont 9 formats JSON standardisés, tous dessinés par un unique composant \code{PlotlyChart.svelte}. Une vue déclare où lire chaque donnée (\texttt{bind}), au lieu de supposer des noms de clés (annexe~\ref{ann:vues}).
\end{enumerate}

\subsubsection{Les règles du générateur}
\begin{itemize}[leftmargin=*, nosep]
  \item \textbf{Un flux aléatoire par variable,} dérivé de la graine et du nom de la variable. Ajouter une variable à un YAML ne change pas les tirages des autres.
  \item \textbf{Des uniformes communes :} toute loi est tirée à partir d'uniformes de ce flux, par sa fonction quantile. Changer de loi déforme les mêmes tirages au lieu d'en produire d'autres : la Référence se compare point par point à l'état cassé, et les courbes de sensibilité sont lisses.
  \item \textbf{Des variables ordonnées par leurs dépendances :} le moteur lit les noms utilisés dans chaque expression et calcule les variables dans le bon ordre ; un cycle est refusé au chargement.
  \item \textbf{Des dimensions déclarées :} une variable vaut par défaut un tableau (répétitions, observations) ; elle peut aussi valoir une valeur par répétition, ou une matrice pour la régression multiple. Les fonctions de tirage ne prennent jamais de taille en argument : la forme vient de la déclaration, ce qui supprime les erreurs de dimensions.
  \item \textbf{Des points ajoutés bien définis :} le générateur indique quelles variables reçoivent les coordonnées d'un point posé au clic (\texttt{inject}). Le point est ajouté à toutes les répétitions, jamais à la Référence.
\end{itemize}

\begin{table}[H]
  \centering
  \small
  \begin{tabularx}{\textwidth}{@{}L{2.6cm} L{5.2cm} Y@{}}
    \toprule
    \textbf{Loi} & \textbf{Tirage depuis une uniforme $u$} & \textbf{Mise à l'échelle} \\
    \midrule
    \texttt{normal} & quantile de la loi normale & moyenne 0, variance 1 \\
    \texttt{uniform} & $\sqrt{3}\,(2u - 1)$ & moyenne 0, variance 1 \\
    \texttt{exponential} & $-\ln(1 - u) - 1$ & moyenne 0, variance 1 \\
    \code{student_2} & $(2u - 1)/\sqrt{2u(1-u)}$ & médiane 0, écart interquartile de la normale (variance infinie) \\
    \texttt{cauchy} & $\tan(\pi(u - 0{,}5))$ & médiane 0, écart interquartile de la normale (pas d'espérance) \\
    \texttt{bimodal} & 2 uniformes : la composante, puis la position & moyenne 0, variance 1 \\
    \bottomrule
  \end{tabularx}
  \caption{Le registre des lois (\texttt{engine/laws.py}). Les lois sont comparées à échelle égale ; chaque entrée indique aussi si l'espérance et la variance existent, pour refuser une mesure de biais sous Cauchy.}
\end{table}

\begin{table}[H]
  \centering
  \small
  \begin{tabularx}{\textwidth}{@{}L{3.6cm} Y L{3.8cm}@{}}
    \toprule
    \textbf{Agrégateur} & \textbf{Calcule} & \textbf{Erreur Monte Carlo} \\
    \midrule
    \code{rate(f)} & part des répétitions où $f$ est vrai (couverture, rejet) & $\sqrt{p(1-p)/n}$ \\
    \code{mean(f)} & moyenne & écart type$/\sqrt{n}$ \\
    \code{bias(f, t)} & moyenne de $f$ moins la vraie valeur $t$ & écart type$/\sqrt{n}$ \\
    \code{rmse(f, t)} & racine de l'erreur quadratique moyenne & non calculée \\
    \code{sd(f)}, \code{var(f)} & dispersion & non calculée \\
    \code{median(f)}, \code{quantile(f, q)} & position & non calculée \\
    \code{first(f)} & valeur sur la répétition 0, c'est-à-dire l'échantillon affiché & non calculée \\
    \code{summary(k)} & grandeur globale renvoyée par la méthode (erreur d'un schéma numérique) & non calculée \\
    \bottomrule
  \end{tabularx}
  \caption{Les agrégateurs (\texttt{engine/aggregators.py}). Ici $n$ est le nombre de répétitions valides. L'argument $f$ peut être une expression sur les champs de \texttt{reps}, par exemple \texttt{rate(reps.pvalue < 0.05)}.}
\end{table}

Une même expérience peut exposer plusieurs hypothèses de types différents (continu, booléen, choix de loi). L'étudiant teste ainsi de vraies interactions, par exemple une loi asymétrique, une hétéroscédasticité modérée et un point levier en même temps.

\textbf{Répétitions invalides.} Un échec numérique sur une répétition (matrice singulière, dépassement) est exclu des mesures et compté ; au-delà de 5\,\% de répétitions invalides, un avertissement s'affiche. Un échec qui est le sujet de l'expérience, comme Euler qui diverge, n'est pas une erreur : la méthode le signale, et il devient une mesure (taux de divergence).

\subsection{Un seul pipeline pour trois disciplines}
Le tableau suivant montre trois théorèmes très différents qui traversent exactement les mêmes étapes, sans toucher ni au moteur ni à l'interface.

\begin{table}[H]
  \centering
  \footnotesize
  \begin{tabularx}{\textwidth}{@{}L{2.5cm} Y Y Y@{}}
    \toprule
    \textbf{Étape} & \textbf{Régression linéaire} & \textbf{Test de Student} & \textbf{Pendule simple (EDO)} \\
    \midrule
    \textbf{1. État} &
    $n=150$, hétéroscédasticité 0,8 &
    $n=25$, loi de Cauchy &
    $\theta_0=1{,}2$ rad, pas $\Delta t = 0{,}01$ s \\
    \midrule
    \textbf{2. Générer} &
    500 échantillons $(X,Y)$ à variance non constante &
    500 échantillons sous $H_0$ et sous $H_1$ (décalage de 0,5) avec la loi choisie &
    Solution de référence de $\ddot{\theta} + \frac{g}{L}\sin\theta = 0$ \\
    \midrule
    \textbf{3. Méthode} &
    MCO : pente, résidus, IC classiques &
    Test de Student : taux de rejet sous $H_0$ et sous $H_1$ &
    Euler explicite \\
    \midrule
    \textbf{4. Correction} &
    Écarts-types HC3 (White) &
    Test de Wilcoxon (rangs) &
    Runge-Kutta d'ordre 4 \\
    \midrule
    \textbf{5. Mesure} &
    Couverture 75\,\% (cible 95\,\%) : Rompu &
    Puissance 13\,\% contre 67\,\% sous loi normale : Rompu &
    Erreur maximale sur l'angle 45\,\% (cible : moins de 5\,\%) : Rompu \\
    \midrule
    \textbf{6. Vues} &
    nuage, résidus, répétitions, sensibilité &
    distribution, répétitions, sensibilité par loi &
    trajectoires $\theta(t)$, sensibilité \\
    \bottomrule
  \end{tabularx}
  \caption{Trois théorèmes, un seul pipeline. Les mesures sont celles obtenues en simulant les expériences avant le hackathon ; elles varient légèrement avec la graine.}
\end{table}

\subsection{Comment l'Entraînement fabrique ses questions}\label{sec:fabrique}
Pour qu'aucune question ne s'apprenne par cœur, le serveur fabrique chaque série à partir d'une graine, en cinq étapes déterministes :
\begin{enumerate}[leftmargin=*, nosep]
  \item \textbf{Tirer les questions et leurs paramètres :} $k$ questions parmi les expériences du filtre, puis une valeur pour chaque paramètre variable (par exemple $n$ entre 150 et 600).
  \item \textbf{Remplir l'énoncé} avec ces valeurs (\og avec $n = 250$\dots \fg{}) et les reporter dans l'état de validation de chaque choix.
  \item \textbf{Vérifier la correction :} lancer la simulation de chaque état de validation et contrôler sa condition \texttt{verify}. Si elle échoue, tirer d'autres paramètres (5 essais au plus) ; sinon, retirer la question de la série et le journaliser.
  \item \textbf{Pour un diagnostic sur graphique,} choisir en secret une pathologie parmi celles fabriquées à partir des hypothèses, lancer \texttt{run()} et n'envoyer que le graphique (sans les paramètres cachés).
  \item \textbf{Mélanger l'ordre des choix} pour qu'aucune position ne soit mémorisable.
\end{enumerate}

\textbf{Des séries prêtes à l'avance.} Fabriquer une série demande entre 30 et 60 exécutions (états de validation et nouveaux tirages), soit plusieurs secondes sur un hébergeur gratuit. L'étudiant ne les attend jamais :
\begin{itemize}[leftmargin=*, nosep]
  \item \textbf{Une réserve :} le serveur garde 3 séries prêtes par filtre (toutes les expériences, et chaque expérience P0). Ouvrir l'Entraînement prend une série de la réserve : la réponse est immédiate.
  \item \textbf{Un remplacement en arrière-plan :} dès qu'une série est prise, le serveur en fabrique une autre, mais seulement quand aucune requête n'est en cours, pour ne jamais ralentir un étudiant ni la démo.
  \item \textbf{Une validation instantanée :} les états de validation sont calculés pendant la fabrication et restent en cache ; l'affichage de la correction après chaque réponse est immédiat.
  \item \textbf{Un lien par série :} une série est entièrement déterminée par son filtre, sa graine et le contenu des YAML. Son lien (\code{?train=regression-lineaire&seed=123&v=1&h=3f09c1}) redonne exactement les mêmes questions. Si la série n'est plus en mémoire (serveur redémarré), elle est refabriquée à l'identique en quelques secondes. Si les YAML ont changé depuis (\texttt{h} différent), l'interface prévient : \og Les questions ont été mises à jour depuis ce lien \fg{}.
  \item \textbf{Une série de démo statique :} \code{scripts/make_fixtures.py} fabrique au build la série de la démo, avec ses résultats de validation, et Cloudflare la sert en JSON statique. Les 15 premières secondes du pitch tournent même si le moteur dort.
\end{itemize}

\subsection{Pourquoi YAML (et JSON seulement entre le serveur et l'interface)}
Les expériences s'écrivent \textbf{uniquement en YAML} ; le JSON ne sert qu'aux échanges entre l'API et l'interface. Il n'y a donc pas de convertisseur.
\begin{itemize}[leftmargin=*, nosep]
  \item \textbf{Les formules :} en YAML, on écrit le LaTeX tel quel ; en JSON, chaque antislash doit être doublé (\texttt{\textbackslash\textbackslash hat\{\textbackslash\textbackslash beta\}}).
  \item \textbf{Les explications :} des paragraphes sur plusieurs lignes et des commentaires \texttt{\#}, naturels en YAML.
  \item \textbf{La relecture :} l'équipe va relire beaucoup de contenu généré par IA ; le YAML se lit plus vite.
\end{itemize}
Les pièges du YAML sont neutralisés par trois règles : lecture par \code{yaml.safe_load}, schéma Pydantic strict qui renvoie le numéro de ligne fautif, et chaînes de texte toujours entre guillemets (sinon YAML peut lire un \texttt{no} comme le booléen faux). Chaque fichier commence par \code{schema_version: 1} ; le moteur refuse une version qu'il ne connaît pas.

\subsection{Un seul espace de noms dans les expressions}
Toutes les expressions du YAML (générateur, arguments des méthodes, défis, conditions \texttt{verify}, contrôles grisés) lisent les mêmes noms :

\begin{table}[H]
  \centering
  \small
  \begin{tabularx}{\textwidth}{@{}L{3.6cm} Y L{4.4cm}@{}}
    \toprule
    \textbf{Nom} & \textbf{Valeur} & \textbf{Utilisable dans} \\
    \midrule
    \code{p.<id>} & valeur d'un paramètre & partout \\
    \code{h.<id>} & valeur d'une hypothèse (nombre, booléen ou clé d'option) & partout \\
    nom nu (\texttt{x}, \texttt{eps}) & variable du générateur & générateur, arguments des méthodes \\
    \code{m.<mesure>} & valeur de la mesure pour la méthode principale & défis, \texttt{verify} \\
    \code{m.<mesure>.<methode>} & valeur pour une méthode donnée & défis, \texttt{verify} \\
    \code{m.<mesure>_se} & erreur Monte Carlo de la mesure & défis, \texttt{verify} \\
    \code{s.<mesure>} & verdict (\texttt{"ok"}, \texttt{"warning"}, \texttt{"broken"}) de la méthode principale & défis, \texttt{verify} \\
    \code{base.<mesure>} & valeur de la Référence pour la méthode principale & défis, \texttt{verify} \\
    \bottomrule
  \end{tabularx}
  \caption{L'espace de noms des expressions. Les paramètres sont toujours préfixés par \texttt{p.}, pour qu'aucun ne se confonde avec une variable du générateur.}
\end{table}

\subsection{Créer une expérience sans payer d'API}
L'écran \emph{Import} \pone{} réunit trois outils :
\begin{enumerate}[leftmargin=*, nosep]
  \item un bouton \textbf{\og Copier le prompt \fg{}} qui fournit le gabarit exact d'une expérience, à coller dans n'importe quel assistant IA externe : aucun appel IA depuis notre app ;
  \item un éditeur avec \textbf{validation ligne par ligne}, les tests scientifiques de la section~\ref{sec:tests} et une prévisualisation immédiate ;
  \item un bouton \textbf{\og Créer un lien de partage \fg{}}.
\end{enumerate}

\subsection{Où tourne l'application}
Crash-Test est une \textbf{web app accessible par un simple lien}. Elle est découpée en trois morceaux, chacun hébergé là où il est le plus à l'aise :

\begin{table}[H]
  \centering
  \small
  \begin{tabularx}{\textwidth}{@{}L{3.3cm} L{4.4cm} Y@{}}
    \toprule
    \textbf{Morceau} & \textbf{Où} & \textbf{Pourquoi} \\
    \midrule
    Interface Svelte & \textbf{Cloudflare}, comme une autre application de l'équipe & Un site statique : gratuit, rapide partout, et c'est ce lien que l'on donne au jury. \\
    Moteur Python (FastAPI) & Un \textbf{hébergeur Python gratuit} qui accepte Docker (Hugging Face Spaces ou Render, choisi avant le hackathon) & Les simulations Monte Carlo ont besoin de NumPy, SciPy et statsmodels natifs, et de plus de temps de calcul que les fonctions Cloudflare gratuites n'en accordent (10 ms par requête). \\
    Liens de partage & \textbf{Supabase} & Une simple table, reprise d'une autre application de l'équipe. \\
    \bottomrule
  \end{tabularx}
  \caption{Les trois morceaux de l'application en ligne.}
\end{table}

Ce découpage ne change presque rien au code : l'adresse de l'API devient une variable d'environnement (\code{VITE_API_URL}) lue par \texttt{executor.ts}, et FastAPI autorise le domaine Cloudflare (CORS). Le moteur est livré dans une image Docker, la même en local et en ligne.

Trois règles :
\begin{itemize}[leftmargin=*, nosep]
  \item \textbf{Mettre en ligne dès le premier soir,} avec la tranche verticale, et non à la fin : c'est là qu'on découvre les surprises (CORS, mémoire, temps de démarrage).
  \item \textbf{Réveiller le serveur avant la soutenance :} un hébergeur gratuit peut mettre le moteur en veille ; on lance une requête 5 minutes avant de passer.
  \item \textbf{Garder une version locale de secours,} servie par FastAPI sur un seul port et sans réseau.
\end{itemize}
Ce même moteur en ligne servira ensuite de \og bloc Python \fg{} à l'application \emph{Concepts \& Révision}.

\subsection{Sécurité, robustesse et temps de calcul}\label{sec:robustesse}
\begin{itemize}[leftmargin=*, nosep]
  \item \textbf{Aucun \texttt{eval()} ni \texttt{exec()} :} toutes les expressions du YAML passent par \texttt{asteval}, avec une liste blanche de fonctions NumPy et des tirages liés à la graine. Pas de boucle ni de compréhension, pas d'attribut commençant par \code{_}, exposants bornés à 100.
  \item \textbf{Aucune valeur aberrante transmise :} \texttt{NaN} et \texttt{Infinity} deviennent \texttt{null}, avec un avertissement en français (vouvoiement). Chaque avertissement a un code stable, pour que l'interface le traite sans lire son texte.
  \item \textbf{Trois règles anti-bugs côté interface :} la graine ne change jamais quand on bouge un curseur (seul \og Nouveau tirage \fg{} l'incrémente) ; les requêtes sont regroupées (100 ms) et les obsolètes annulées pendant un glissement ; la caméra 3D ne se réinitialise jamais (\code{Plotly.react()} avec \code{uirevision: 'keep'}).
  \item \textbf{Des calculs qui ne se bloquent pas :} les routes sont écrites en \texttt{def} et non en \texttt{async def}, pour que FastAPI exécute chaque calcul dans un fil séparé : un balayage long ne bloque pas les autres requêtes. Un seul processus suffit pour l'équipe et la démo.
  \item \textbf{Un cache à chaque étape :} le monde généré, la Référence et les balayages sont mis en cache. La clé combine la version du format, le contenu exact du YAML et l'état de calcul mis sous forme canonique (valeurs par défaut complétées, valeurs arrondies au pas du contrôle, clés triées). Deux états équivalents partagent donc le même résultat, et les états de la démo sont précalculés.
  \item \textbf{Un budget de calcul fixe} \pone{} : une unité de coût est une observation traitée par une méthode (répétitions $\times$ $n$). Le budget vaut 4 millions d'unités pour \code{/api/run} et 30 millions pour \code{/api/sweep}, ce qui correspond à la limite de 200 ms par méthode de l'annexe~\ref{ann:methodes}. Au-delà, le moteur divise le nombre de répétitions jusqu'à tenir le budget, sans descendre sous 100, et le signale (\og Estimation sur 200 tirages \fg{}). Sous 100, il refuse avec un message clair. Si l'hébergeur est plus lent que prévu, rien n'est à recalibrer : le moteur réduit et prévient. En \pzero{}, les plafonds et les délais du tableau ci-dessous suffisent.
\end{itemize}

Les mêmes plafonds s'appliquent aux expériences officielles et importées. Un YAML importé s'exécute sur notre serveur : sans plafonds, un seul lien partagé suffirait à saturer l'hébergeur gratuit.

\begin{table}[H]
  \centering
  \small
  \begin{tabularx}{\textwidth}{@{}Y L{4.2cm} L{4.4cm}@{}}
    \toprule
    \textbf{Limite} & \textbf{Valeur} & \textbf{Vérifiée} \\
    \midrule
    Répétitions (\texttt{repeat}) & 1\,000 au plus & au chargement \\
    Taille d'échantillon (bornes de $n$) & 5\,000 au plus & au chargement \\
    Méthodes par expérience & 4 au plus & au chargement \\
    Variables du générateur & 30 au plus & au chargement \\
    Longueur d'une expression & 300 caractères au plus & au chargement \\
    Taille du fichier YAML & 64 Ko au plus & au chargement et à l'envoi \\
    Points ajoutés à la main & 20 au plus & à chaque requête \\
    Grille d'un balayage & 21 points, ou 144 cases & à chaque requête \\
    Délai de \code{/api/run} & 2 s & à l'exécution \\
    Délai de \code{/api/sweep} & 15 s & à l'exécution \\
    Fabrication d'une série d'Entraînement à la demande (lien partagé hors réserve) & 15 s & à l'exécution \\
    Réserve de séries prêtes & 3 par filtre & au démarrage du serveur \\
    \bottomrule
  \end{tabularx}
  \caption{Les plafonds de calcul, identiques pour toutes les expériences.}
\end{table}

\subsection{Les tests qui garantissent chaque expérience}\label{sec:tests}
Le fichier \code{tests/test_experiments.py} applique dix vérifications à chaque YAML du dossier \code{experiments/}, et à tout YAML soumis à l'écran Import (avec moins de répétitions). Les tests 1, 3, 4, 8 et 10 sont \pzero{} ; les autres sont \pone{}.
\begin{enumerate}[leftmargin=*, nosep]
  \item \textbf{Le théorème tient à l'état sain :} la mesure primaire de la méthode principale est Fiable.
  \item \textbf{Chaque hypothèse sert à quelque chose :} cassée seule au maximum (ou sur chaque option), elle fait sortir au moins une mesure de Fiable, sauf si elle est marquée \code{robust_expected} (la loi uniforme pour un test de Student, par exemple).
  \item \textbf{Chaque correction répare ce qu'elle annonce :} une correction ou une alternative déclare les hypothèses qu'elle répare (\texttt{repairs}) ; sur chaque état où la méthode principale est Rompue à cause de l'une d'elles, la méthode qui la répare revient à Fiable ou Fragile. HC3 répare l'hétéroscédasticité, pas l'autocorrélation : c'est écrit dans le YAML et vérifié.
  \item \textbf{Les contre-exemples font ce qu'ils annoncent :} chacun rompt au moins une mesure, sauf ceux marqués \code{expect: same_summary} (Anscombe), pour lesquels le test vérifie au contraire que pente et $R^2$ sont identiques à $10^{-9}$ près.
  \item \textbf{Les défis sont atteignables :} un balayage sur leurs contrôles trouve au moins un état qui les valide.
  \item \textbf{Les questions tiennent :} la condition \texttt{verify} passe sur au moins 90\,\% de 50 tirages de paramètres.
  \item \textbf{Les pathologies se distinguent :} deux pathologies d'une même question diffèrent d'au moins 3 erreurs Monte Carlo sur une mesure.
  \item \textbf{Le calcul est déterministe :} deux exécutions du même état donnent exactement le même résultat.
  \item \textbf{Le budget est tenu :} l'état par défaut répond en moins de 500 ms en local.
  \item \textbf{Les explications sont complètes :} chaque hypothèse a \texttt{why} et \code{when_broken}, chaque correction et chaque alternative ont \texttt{explain}.
\end{enumerate}
En plus, des tests par propriétés (bibliothèque \texttt{hypothesis}) tirent 200 états au hasard dans les bornes de chaque expérience : aucune exception, aucun \texttt{NaN} ni \texttt{Infinity} dans la réponse. Pour que le test 8 passe aussi sur l'hébergeur, le calcul matriciel tourne sur un seul fil (\code{OMP_NUM_THREADS=1}) et les versions de NumPy et SciPy sont fixées dans \code{uv.lock}.

\subsection{Quand un besoin imprévu apparaît}\label{sec:imprevu}
Le format YAML, l'API et le contrat des méthodes décrits en annexe restent les mêmes pendant les 30 heures. Un besoin imprévu passe par une des sorties prévues ci-dessous, jamais par une modification de ces contrats :

\begin{table}[H]
  \centering
  \small
  \begin{tabularx}{\textwidth}{@{}L{5.6cm} Y@{}}
    \toprule
    \textbf{Besoin} & \textbf{Sortie prévue} \\
    \midrule
    Une grandeur qu'aucun agrégateur ne calcule & La méthode la renvoie dans \texttt{summary}, lue par \code{summary(k)}. \\
    Un agrégateur manquant & Une fonction de plus dans \code{engine/aggregators.py}, sans toucher au format. \\
    Une loi absente du registre & Une option avec \texttt{expr:} (expression libre) à la place de \texttt{law:} ; ses tirages ne sont pas couplés aux autres lois. \\
    Une nouvelle méthode & Une fonction de plus avec \code{@register_method}. \\
    Un réglage d'affichage propre à une vue & Le champ libre \texttt{ui: \{\}} de la vue, ignoré par le moteur et transmis tel quel à l'interface. \\
    Une série de plus dans une vue & Une seconde vue du même type. \\
    Une vue qui ne se construit pas & La vue de secours \code{measures_table} s'affiche à sa place. \\
    \bottomrule
  \end{tabularx}
  \caption{Les sorties prévues pour les besoins imprévus.}
\end{table}

% ============================================================================
% SECTION 5 : CATALOGUE
% ============================================================================
\section{Le catalogue d'expériences}

\begin{enbref}
Trois expériences au cœur du hackathon (régression, test de Student, pendule et méthode d'Euler), une en extension si le P0 est vert à H+21 (LGN et TCL), et une réserve pour la suite. Pour chacune : 2 défis et 3 à 5 QCM solides plutôt qu'une grande banque. Les chiffres de ce cahier ont été vérifiés par simulation avant le hackathon.
\end{enbref}

\begin{table}[H]
  \centering
  \footnotesize
  \begin{tabularx}{\textwidth}{@{}L{2.6cm} c Y Y Y@{}}
    \toprule
    \textbf{Expérience} & \textbf{Prio.} & \textbf{Hypothèses cassées} & \textbf{Méthodes} & \textbf{Ce qu'on montre} \\
    \midrule
    \textbf{1. Régression linéaire} & \pzero &
    \pzero{} variance constante (curseur), loi des erreurs (choix), point levier (interrupteur). \pone{} indépendance (erreurs et $X$ en AR(1)), forme de la relation (choix), préréglages Anscombe. &
    MCO classique / écarts-types HC3 ; \pone{} Newey-West en alternative. &
    Couverture qui passe de 94\,\% à 75\,\% quand la variance dépend de $x$, HC3 qui la ramène vers 94\,\% ; défi \og piéger Breusch-Pagan \fg{} ; la loi des erreurs ne change presque rien (résultat à retenir). \textbf{Expérience de la démo.} \\
    \midrule
    \textbf{2. Test de Student} & \pzero &
    Loi sous $H_0$, toutes symétriques : Normale, Uniforme, Bimodale, Student à 2 ddl, Cauchy ; petit ou grand $n$. $H_0$ porte sur le centre de symétrie. &
    Test de Student / test de Wilcoxon. &
    Sous des queues lourdes, le test de Student devient trop prudent (risque réel 3\,\% au lieu de 5\,\%) et sa puissance s'effondre (13\,\% contre 67\,\% à $n = 25$ sous Cauchy) ; Wilcoxon garde son niveau et l'essentiel de sa puissance. \\
    \midrule
    \textbf{3. Pendule et méthode d'Euler (EDO)} & \pzero &
    Pas de temps $\Delta t$ (curseur, de 0,0003 à 0,1 s) ; amplitude $\theta_0$ en paramètre. &
    Euler explicite / Runge-Kutta d'ordre 4. &
    Erreur qui grandit avec le pas, puis la falaise : le pendule fait le tour de son axe. RK4 reste exact avec un pas 100 fois plus grand. \textbf{Preuve que ce n'est pas qu'une app de stats.} \\
    \midrule
    \textbf{4. LGN et TCL} & \pone &
    Loi (Normale, Exponentielle, Student à 2 ddl, Cauchy sans espérance), $n$ croissant. &
    Moyenne empirique et intervalle de Student / médiane empirique. &
    Faisceau de 30 trajectoires ; intervalle de Student sous loi asymétrique (couverture 90\,\% à $n = 10$, réparée par le grand $n$) ; Cauchy, que rien ne répare. \\
    \midrule
    \textbf{5. Régression multiple 3D} & \ptwo &
    Corrélation entre $X_1$ et $X_2$ (multicolinéarité), point levier en 3D. &
    MCO / régression Ridge. &
    Plan ajusté en 3D, compromis biais-variance. \\
    \midrule
    \textbf{6. Continuité et TFA} & \ptwo &
    Fonction libre $f(x)$ et préréglages, saut de discontinuité. &
    Dérivée de $F$ / découpage de Chasles. &
    Saisie libre de la fonction du jury en direct. \\
    \midrule
    \textbf{7. Pendule et petits angles} & \ptwo &
    Amplitude $\theta_0$ (approximation $\sin\theta \approx \theta$). &
    Modèle linéaire / formule de Borda. &
    Erreur de période qui grandit avec l'amplitude. \\
    \bottomrule
  \end{tabularx}
  \caption{Le catalogue : socle P0, extension P1 et réserve P2.}
\end{table}

\begin{rulebox}{Une erreur évitée grâce à la simulation}
La version précédente du cahier corrigeait le test de Student \emph{sur une loi exponentielle} par le test de Wilcoxon. C'est faux : Wilcoxon teste le centre de symétrie, pas la moyenne. Sous une loi exponentielle centrée sur sa moyenne, il rejette à tort dans 15\,\% des cas à $n = 25$ et dans 39\,\% des cas à $n = 100$. D'où deux décisions : le test de Student n'utilise que des lois symétriques (où Wilcoxon est une vraie correction), et l'asymétrie est traitée dans LGN et TCL, où c'est le grand échantillon qui répare. La règle qui en découle : toute paire méthode/correction est simulée par R2 avant d'être écrite dans un YAML.
\end{rulebox}

\textbf{Réserve} \ptwo{} : moteur \texttt{matrix} et vue espace vectoriel 2D/3D (conditionnement, bases, SVD), moteur \texttt{process} (séries temporelles), miniatures de sensibilité sur les cartes de la galerie. Les cours de l'ENSAI et de Dauphine offrent des dizaines d'autres candidats (Gauss-Markov et biais de variable omise, biais-variance en kNN, descente de gradient non convexe, tâtonnement walrasien, données manquantes MNAR\dots).

% ============================================================================
% SECTION 6 : ORGANISATION
% ============================================================================
\section{Organisation de l'équipe}

\begin{enbref}
Deux binômes \emph{mixtes} se relaient : le binôme de jour (Kilian et R2 : interface, méthodes et contenu) et le binôme de nuit (R1 et R4 : moteur et écrans secondaires). L'interface et les méthodes avancent dès H+4 sur de faux résultats JSON, sans attendre le moteur. À chaque relève, le binôme qui arrive relit et fusionne le travail de celui qui part. Tout le monde code avec l'IA : le goulot devient la relecture et l'intégration, d'où des règles strictes.
\end{enbref}

\subsection{Les 4 rôles}
\begin{table}[H]
  \centering
  \small
  \begin{tabularx}{\textwidth}{@{}L{3.4cm} Y L{3.6cm}@{}}
    \toprule
    \textbf{Rôle} & \textbf{Responsabilités} & \textbf{Relit les changements de\dots} \\
    \midrule
    \textbf{R1 : Moteur et API} & Orchestrateur \texttt{run()} et balayage, générateurs, flux aléatoires, cache et budget de calcul, routes, génération procédurale des QCM, fixtures pour l'interface, tests du moteur et de l'API. & \code{backend/engine}, \code{backend/api} \\
    \midrule
    \textbf{R2 : Maths et contenu} & Méthodes et corrections (simulées avant d'être écrites), registre des lois, agrégateurs, cibles et tolérances des mesures, les 3 fichiers YAML P0 puis LGN et TCL (justesse mathématique, défis, QCM, explications), tests des expériences. & \code{backend/methods}, \code{experiments/} \\
    \midrule
    \textbf{R3 (Kilian) : Architecte interface et intégrateur} & Squelette Svelte, \code{PlotlyChart}, client API, types générés, rôles et styles, équation dynamique et sa règle de couleur, Laboratoire, Entraînement, panneau Comprendre, lien de partage. Seul de l'équipe à avoir déjà fait du Svelte. & \code{frontend/} \\
    \midrule
    \textbf{R4 : Écrans secondaires et qualité} & Accueil et galerie, écran Bilan et profil démo, écran Import, test du chargeur, premiers jets des explications (relus par R2), audit du vouvoiement, tests utilisateurs, slides et vidéo de secours. & ses propres écrans, relus par R3 \\
    \bottomrule
  \end{tabularx}
  \caption{Répartition des rôles.}
\end{table}

Les tâches de R4 sont des écrans indépendants, qui ne bloquent personne et se vérifient à l'\oe il : avec l'IA et les faux JSON, elles avancent sans dépendre du reste. De même, les méthodes de R2 sont des fonctions pures, testables seules avec \texttt{pytest} avant même que le moteur existe.

\subsection{Le planning en relais}
\begin{itemize}[nosep]
  \item \textbf{Binôme de jour :} Kilian (R3) + R2, l'interface principale, les méthodes et le contenu des expériences.
  \item \textbf{Binôme de nuit :} R1 + R4, le moteur, l'API et les écrans secondaires.
\end{itemize}
Le hackathon dure 30 heures. Les heures indiquées supposent un départ à 9\,h (H0) ; elles se décalent si le hackathon commence à une autre heure. Chaque binôme dort environ 7 heures (le binôme de jour de 21\,h à 4\,h, le binôme de nuit de 6\,h à 11\,h).

\begin{figure}[H]
  \centering
  \begin{tikzpicture}[x=0.43cm, y=1cm]
    % Libelles des lignes
    \node[anchor=east, font=\small\bfseries, text=ctMethod] at (-0.5,0) {Jour};
    \node[anchor=east, font=\small\bfseries, text=ctGenerator] at (-0.5,-1.1) {Nuit};

    % Creneaux communs (tous les 4)
    \foreach \a/\b in {0/4, 11/12, 19/21, 26/30}{
      \fill[ctPrimary!75, rounded corners=2pt] (\a,0.35) rectangle (\b,-1.45);
    }
    % Relais binome de jour
    \foreach \a/\b in {4/11, 21/26}{
      \fill[ctMethod!80, rounded corners=2pt] (\a,0.35) rectangle (\b,-0.35);
    }
    % Relais binome de nuit
    \fill[ctGenerator!80, rounded corners=2pt] (12,-0.75) rectangle (19,-1.45);

    % Textes dans les barres
    \node[font=\footnotesize\bfseries, text=white] at (2,-0.55) {Tous};
    \node[font=\footnotesize\bfseries, text=white] at (7.5,0) {Kilian + R2};
    \node[font=\footnotesize\bfseries, text=white] at (15.5,-1.1) {R1 + R4};
    \node[font=\footnotesize\bfseries, text=white, align=center] at (20,-0.55) {P0};
    \node[font=\footnotesize\bfseries, text=white] at (23.5,0) {Kilian + R2};
    \node[font=\footnotesize\bfseries, text=white] at (28,-0.55) {Tous};

    % Axe des heures
    \draw[ctBaseline] (0,-1.75) -- (30,-1.75);
    \foreach \h in {0,4,11,12,19,21,26,28,30}{
      \draw[ctBaseline] (\h,-1.7) -- (\h,-1.8);
    }
    \foreach \h/\lab/\clk in {0/H0/9h, 4/H4/13h, 11/H11/20h, 19/H19/4h, 21/H21/6h, 26/H26/11h, 30/H30/15h}{
      \node[font=\scriptsize, text=ctBaseline, align=center] at (\h,-2.15) {\lab\\[-1pt]\clk};
    }

    % Jalons
    \foreach \h/\lab in {4/{Tranche verticale}, 21/{Jalon P0}, 28/{Code final}}{
      \node[diamond, fill=ctBroken, inner sep=1.6pt] at (\h,-2.7) {};
      \node[font=\scriptsize\bfseries, text=ctBroken] at (\h,-3.05) {\lab};
    }
  \end{tikzpicture}
  \caption{Les relais sur 30 heures (départ supposé à 9\,h). En bleu foncé, les créneaux où les 4 membres sont présents.}
\end{figure}

\renewcommand{\arraystretch}{1.15}
\begin{longtable}{@{}L{1.9cm} L{2.4cm} >{\raggedright\arraybackslash}p{\dimexpr\textwidth-4.3cm-4\tabcolsep\relax}@{}}
  \toprule
  \textbf{Créneau} & \textbf{Présents} & \textbf{Objectifs et livrables} \\
  \midrule
  \endfirsthead
  \toprule
  \textbf{Créneau} & \textbf{Présents} & \textbf{Objectifs et livrables} \\
  \midrule
  \endhead
  \bottomrule
  \endfoot

  \textbf{H0 à H4} \newline {\small 9\,h à 13\,h} &
  \textbf{Tous} \newline \emph{Démarrage} &
  Les 4 premières heures servent à coder les contrats décrits en annexe, pas à les discuter.
  \begin{itemize}[leftmargin=*, nosep]
    \item \textbf{H0 à H1 (les 4, R2 au clavier) :} dépôt, \texttt{CLAUDE.md}, puis le fichier \code{regression_lineaire.yaml} recopié \emph{ensemble} depuis l'annexe~\ref{ann:yaml} : il sert de modèle aux autres.
    \item \textbf{H1 à H2 :} \code{engine/schema.py}, qui charge ce fichier et celui du pendule (R1) ; \code{method_contract.py}, \code{aggregators.py} et \code{laws.py} avec leur test de conformité (R2).
    \item \textbf{H2 à H3 :} \code{contracts.py} (requêtes et réponses de \code{/api/run}, \code{/api/sweep} et \code{/api/training/session}, R1) ; \code{roles.ts}, \texttt{tokens.css} et la règle de couleur de \code{equation.ts} (Kilian).
    \item \textbf{H3 à H4 :} \code{scripts/gen_types.sh} et le premier \code{types.gen.ts} (Kilian) ; faux JSON des vues P0 et d'une session d'Entraînement, validés par \code{contracts.py} (R1) ; test du chargeur, qui accepte les YAML de référence et refuse trois YAML volontairement faux avec le bon numéro de ligne (R4).
    \item Installation FastAPI + Vite/Svelte 5 + proxy \code{/api} ; projet Cloudflare relié au dépôt (Kilian), \texttt{Dockerfile} du moteur sur l'hébergeur choisi avant le hackathon (R1).
    \item \textbf{Jalon tranche verticale (R1) :} une première version de \code{POST /api/run} (générateur \texttt{random} et MCO) renvoie une vraie courbe affichée par \code{PlotlyChart.svelte}. On pose alors le tag git \code{contracts-v1}.
  \end{itemize} \\
  \midrule

  \textbf{H4 à H11} \newline {\small 13\,h à 20\,h} &
  \textbf{Jour} \newline Kilian + R2 \newline \emph{R1 et R4 se reposent} &
  \begin{itemize}[leftmargin=*, nosep]
    \item \textbf{Kilian :} squelette Svelte, client API, \code{PlotlyChart} (rôles, styles, axes stables), équation dynamique, puis les deux c\oe urs sur faux JSON : le \textbf{Laboratoire} (3 types de contrôles, onglets, Correction, Référence, verdicts) et un \textbf{Entraînement minimal} (une question, les 3 crans, la validation par simulation). Interface en ligne sur Cloudflare. \textbf{Avant la relève :} un composant modèle et une fiche par écran de la nuit (voir ci-dessous).
    \item \textbf{R2 :} méthodes P0 en fonctions pures, vectorisées et testées seules contre \texttt{statsmodels} et SciPy (MCO, HC3 et Breusch-Pagan ; Student et Wilcoxon ; Euler et RK4), registre des lois, puis les 3 YAML P0 (paramètres, mesures et tolérances, défis, premières questions avec \texttt{verify}).
  \end{itemize} \\
  \midrule

  \textbf{H11 à H12} \newline {\small 20\,h à 21\,h} &
  \textbf{Tous} \newline \emph{Relève du soir} &
  Passation, fusion sur \texttt{main}, revue des vues et des méthodes avec R1, priorités de la nuit. \textbf{Première mise en ligne complète :} l'interface Cloudflare appelle la tranche verticale du moteur hébergé, vérifiée depuis un téléphone. \\
  \midrule

  \textbf{H12 à H19} \newline {\small 21\,h à 4\,h} &
  \textbf{Nuit} \newline R1 + R4 \newline \emph{Kilian et R2 dorment} &
  \begin{itemize}[leftmargin=*, nosep]
    \item \textbf{R1 :} \texttt{expr.py}, \code{rng.py}, \texttt{sanitize.py}, \texttt{runner.py} (exécution, Référence, balayage à 1 axe, défis, cache), générateurs \texttt{random} et \texttt{ode}, branchement des méthodes et des YAML de R2, routes P0 dont la fabrication des séries avec \texttt{verify} et leur réserve en arrière-plan, \code{scripts/make_fixtures.py} pour remplacer les faux JSON par de vrais résultats, tests P0 au vert.
    \item \textbf{R4 :} accueil et galerie, écran Bilan (\texttt{calibration.ts}, profil démo), premiers jets des explications (relus par R2 au matin), liste de tests utilisateurs sur le Laboratoire.
  \end{itemize} \\
  \midrule

  \textbf{H19 à H21} \newline {\small 4\,h à 6\,h} &
  \textbf{Tous} \newline \emph{Jalon P0} &
  Relecture et fusion du travail de nuit, interface branchée sur la vraie API. \textbf{Critère du jalon :} sur les 3 expériences P0, le Laboratoire (casser, Correction, Référence, verdicts) et l'Entraînement minimal (question générée et vérifiée, crans, validation par simulation) tournent sur la vraie API, et les tests P0 de la section~\ref{sec:tests} passent. Sinon, \textbf{application de la liste de coupes}. Choix des extensions P1, dans l'ordre de la section~\ref{sec:perimetre}. \\
  \midrule

  \textbf{H21 à H26} \newline {\small 6\,h à 11\,h} &
  \textbf{Jour} \newline Kilian + R2 \newline \emph{R1 et R4 dorment} &
  \begin{itemize}[leftmargin=*, nosep]
    \item \textbf{Kilian :} Entraînement complet (séries de questions, diagnostics sur graphique, bouton vers le Laboratoire), panneau Comprendre, onglet Défis, lien de partage, puis les extensions P1 retenues côté interface (banc de méthodes, Domaine de validité, balayage automatique).
    \item \textbf{R2 :} explications et QCM définitifs (3 à 5 par expérience, avec leurs conditions \texttt{verify}), réglage des tolérances, précalcul et mise en cache des états de la démo, puis les extensions P1 retenues côté contenu (LGN et TCL, hypothèses P1 de la régression).
  \end{itemize} \\
  \midrule

  \textbf{H26 à H28} \newline {\small 11\,h à 13\,h} &
  \textbf{Tous} \newline \emph{Sprint final} &
  \textbf{Démo complète de bout en bout} (Entraînement, validation par simulation, Laboratoire, Comprendre, Défis, Correction, Bilan, lien de partage), intégration des extensions P1 terminées, corrections issues des tests utilisateurs, dernière passe sur les contenus. \\
  \midrule

  \textbf{H28 à H30} \newline {\small 13\,h à 15\,h} &
  \textbf{Tous} \newline \emph{Finitions et pitch} &
  \begin{itemize}[leftmargin=*, nosep]
    \item \textbf{H28 : fin du développement,} plus aucune fonctionnalité, seulement des corrections de bugs.
    \item Audit : vouvoiement, zéro erreur console, liens de partage, version en ligne testée sur ordinateur et sur téléphone, version locale de secours (un seul port, hors réseau).
    \item Enregistrement de la vidéo de secours, répétitions chronométrées du pitch.
  \end{itemize} \\
\end{longtable}
\renewcommand{\arraystretch}{1.0}

\subsection{Travailler à quatre avec l'IA}
\begin{keybox}{Les règles du jeu}
  \begin{enumerate}[leftmargin=*, nosep]
    \item \textbf{Un dossier par personne.} Le \texttt{CLAUDE.md} précise à chaque assistant IA les dossiers qu'il a le droit de modifier. Les fichiers de contrat (\texttt{schema.py}, \texttt{contracts.py}, \code{method_contract.py}) ne changent pas pendant le hackathon : un besoin imprévu passe par les sorties prévues (section~\ref{sec:imprevu}). \code{types.gen.ts} est produit par un script et ne se modifie jamais à la main.
    \item \textbf{De petits changements fréquents} plutôt qu'une grosse fusion en fin de relais : l'IA produit vite, mais les erreurs s'accumulent en silence.
    \item \textbf{Les tests comme garde-fou :} rien n'entre dans \texttt{main} sans \texttt{pytest -q}, \texttt{ruff check} et \texttt{npm run check} au vert. Un test échoue si \code{types.gen.ts} ne correspond plus à \code{contracts.py}. Le test des expériences applique les dix vérifications de la section~\ref{sec:tests}.
    \item \textbf{Un humain relit chaque changement} avant de le fusionner, même rapidement : Kilian pour l'interface, R1 pour le moteur, R2 pour les YAML et les méthodes.
    \item \textbf{La nuit, l'interface avance sur des branches.} Les écrans produits la nuit sont relus par Kilian et fusionnés au jalon du matin ; les premiers jets d'explications sont relus par R2. Le moteur fait exception : R1 fusionne ses propres changements dès que les tests sont au vert, pour que l'API soit prête au réveil. Avant la relève du soir, le binôme de jour fusionne son propre travail.
    \item \textbf{Une note de passation en 5 lignes} à chaque relève : ce qui tourne sur \texttt{main}, ce qui est en cours, ce qui bloque, la priorité du binôme suivant.
  \end{enumerate}
\end{keybox}

\subsection{Préparer le travail de la nuit sur l'interface}
Kilian est le seul de l'équipe à avoir déjà utilisé Svelte. Pour que les écrans de la nuit (accueil et galerie, Bilan, Import) soient cohérents avec le reste et rapides à relire au matin, il prépare avant la relève de 21\,h :
\begin{itemize}[leftmargin=*, nosep]
  \item \textbf{un composant modèle} (\code{frontend/src/components/_Modele.svelte}) : la syntaxe Svelte 5 attendue (\texttt{\$props}, \texttt{\$state}, \texttt{\$derived}), l'import de \texttt{tokens.css}, le chargement d'une fixture, la mise en page type et le vouvoiement ;
  \item \textbf{une fiche par écran} : une maquette rapide (capture ou croquis), les données reçues (types exacts dans \code{types.gen.ts}), la fixture à utiliser, et une courte liste de vérifications (vouvoiement, aucune erreur console, \texttt{npm run check} au vert) ;
  \item \textbf{les règles Svelte 5 dans le \texttt{CLAUDE.md}}, pour que les assistants IA n'écrivent pas de Svelte 4 (pas de \texttt{export let}, pas de \texttt{\$:}).
\end{itemize}
Ces écrans sont développés sur des branches et fusionnés au jalon du matin, après relecture. Dès que le moteur répond, \code{scripts/make_fixtures.py} régénère les fixtures à partir de vrais résultats : les écrans de la nuit travaillent sur de vraies sorties.

\subsection{Le périmètre en 30 heures}\label{sec:perimetre}
Avec 30 heures au lieu de 36, le P0 est resserré à ce qui fait la démo et la promesse du projet. Tout le reste est \pone{} et n'est commencé qu'après le jalon de H21, dans l'ordre de la liste ci-dessous. Les contrats, eux, couvrent déjà le P1 et le P2 : rien ne change dans le format quand on ajoute une extension.

\begin{table}[H]
  \centering
  \small
  \begin{tabularx}{\textwidth}{@{}L{2.4cm} Y Y Y@{}}
    \toprule
    & \textbf{\pzero{} : le socle (jusqu'à H21)} & \textbf{\pone{} : si le P0 est vert} & \textbf{\ptwo{} : après le hackathon} \\
    \midrule
    \textbf{Expériences} & Régression (variance, loi des erreurs, levier), test de Student, pendule et Euler & LGN et TCL ; régression : indépendance et Newey-West, forme de la relation et Anscombe & Régression 3D, TFA, pendule et petits angles \\
    \midrule
    \textbf{Moteur} & \texttt{run()}, générateurs \texttt{random} et \texttt{ode}, méthodes en liste, agrégateurs, verdicts, Référence, balayage à 1 axe, cache, plafonds et délais & Balayage à 2 axes, budget avec réduction automatique, \code{/api/validate} & Générateur \texttt{function}, \texttt{process}, \texttt{matrix} \\
    \midrule
    \textbf{Laboratoire} & Comprendre, 3 types de contrôles, Correction, Référence, verdicts, équation dynamique, Répétitions, Sensibilité, contre-exemples, 2 défis & Banc de méthodes, Domaine de validité, balayage automatique, ajout de points, journal des tirages & Inspecteur, calques pas à pas \\
    \midrule
    \textbf{Entraînement} & Idées fausses et diagnostics, \texttt{verify}, crans, validation par simulation, pont vers le Laboratoire, Bilan & Questions de prédiction & Révisions espacées \\
    \midrule
    \textbf{Partage et qualité} & Lien de partage par URL, tests 1, 3, 4, 8, 10, types générés, fixtures & Écran Import et partage Supabase, tests 2, 5, 6, 7, 9, tests par propriétés & Comptes facultatifs \\
    \bottomrule
  \end{tabularx}
  \caption{Le périmètre en 30 heures.}
\end{table}

\subsection{La liste de coupes, décidée à l'avance}
Les extensions P1 se font dans l'ordre inverse de cette liste. Si le P0 n'est pas vert au jalon de H+21, on coupe \textbf{dans cet ordre, sans débat} :
\begin{enumerate}[leftmargin=*, nosep]
  \item le journal des tirages et l'ajout de points par clic ;
  \item l'écran Import et le partage Supabase ;
  \item les questions de prédiction ;
  \item les hypothèses P1 de la régression (indépendance, forme de la relation, Anscombe) ;
  \item le balayage automatique ;
  \item le Domaine de validité ;
  \item l'expérience LGN et TCL ;
  \item le banc de méthodes (la Correction seule reste).
\end{enumerate}

\begin{okbox}{Ce qu'on ne coupe jamais}
Le Laboratoire avec Correction, Référence et verdicts, l'Entraînement avec questions vérifiées et validation par simulation, les explications, les vues Répétitions et Sensibilité, l'écran Bilan, le lien de partage, et les 3 expériences P0 au complet.
\end{okbox}

% ============================================================================
% SECTION 7 : DEMO ET PITCH
% ============================================================================
\section{La démo et le pitch}

\begin{enbref}
Le jury doit voir un graphique avant de lire une ligne de texte. On n'ouvre donc pas sur l'accueil : un lien ouvre directement la question de démonstration, avec une graine fixée pour que le résultat soit identique à chaque passage.
\end{enbref}

\subsection{Les 15 premières secondes}
À répéter jusqu'à pouvoir les faire les yeux fermés :
\begin{enumerate}[leftmargin=*, nosep]
  \item L'écran s'ouvre sur un graphique de résidus en entonnoir et la question \og Quelle situation observez-vous ? \fg{}. Aucun texte d'introduction.
  \item Un juge choisit une réponse et un niveau de confiance (ou l'équipe le fait pour lui).
  \item La validation par simulation lance 100 intervalles de confiance : les ratés s'allument en rouge et la couverture tombe autour de 75\,\%, verdict Rompu.
  \item Clic sur Correction : les intervalles HC3 se recalculent en vert et la couverture remonte vers 94\,\%.
\end{enumerate}
Les valeurs exactes seront celles produites par la graine de démo retenue.

\subsection{Le déroulé sur 5 minutes}
\begin{table}[H]
  \centering
  \small
  \begin{tabularx}{\textwidth}{@{}l Y Y@{}}
    \toprule
    \textbf{Temps} & \textbf{Ce qu'on montre} & \textbf{Le message} \\
    \midrule
    0:00 à 0:15 & Le script ci-dessus & \og Ici, la correction est une expérience. \fg{} \\
    0:15 à 1:00 & Le constat : réciter des hypothèses, des QCM corrigés en une ligne & Un problème que tout le jury a vécu \\
    1:00 à 2:00 & \og Ouvrir dans le Laboratoire \fg{}, panneau Comprendre, défi Breusch-Pagan réussi en direct & On comprend, on manipule, on découvre un angle mort \\
    2:00 à 2:45 & L'écran Bilan avec le profil démo & On mesure l'illusion de compétence \\
    2:45 à 3:30 & Le pendule qui fait le tour de son axe avec Euler, et reste exact avec RK4 & La même mécanique hors des statistiques \\
    3:30 à 4:15 & Un fichier YAML, un lien de partage, et l'écran Import s'il est prêt & Ajouter un théorème, c'est écrire un fichier \\
    4:15 à 5:00 & La suite : l'intégration à Concepts \& Révision, d'autres disciplines & Le projet dépasse le hackathon \\
    \bottomrule
  \end{tabularx}
  \caption{Proposition de déroulé, à ajuster au format exact de la soutenance.}
\end{table}

\subsection{Le plan de secours}
\begin{itemize}[leftmargin=*, nosep]
  \item L'application tourne en ligne (lien Cloudflare), avec le moteur réveillé 5 minutes avant de passer ; en secours, elle tourne en local, sans réseau, sur un seul port.
  \item Les états de la démo sont précalculés et en cache ; la série d'Entraînement de la démo est un fichier statique servi par Cloudflare.
  \item Une vidéo de 90 secondes du parcours complet est prête sur deux ordinateurs.
\end{itemize}

% ============================================================================
% SECTION 8 : RISQUES ET SUITE
% ============================================================================
\section{Risques et suite du projet}

\begin{table}[H]
  \centering
  \small
  \begin{tabularx}{\textwidth}{@{}L{5.2cm} Y@{}}
    \toprule
    \textbf{Risque} & \textbf{Garde-fou} \\
    \midrule
    Trop de fonctionnalités, aucune finie & Périmètre en 30 heures (section~\ref{sec:perimetre}), liste de coupes appliquée à H+21 \\
    Une correction fausse mathématiquement & Chaque paire méthode/correction simulée avant d'être écrite ; le cas Wilcoxon sous loi asymétrique a été retiré pour cette raison \\
    Ressembler à une énième app de QCM & Premier écran toujours graphique, démo qui s'ouvre sur un graphique \\
    Une erreur mathématique dans une expérience & Les dix tests de la section~\ref{sec:tests} : état sain Fiable, chaque hypothèse utile, Correction qui répare, questions vérifiées \\
    Une question dont la correction est fausse & Condition \texttt{verify} contrôlée à chaque génération et testée sur 50 tirages \\
    L'interface et le moteur ne parlent plus le même format & Types TypeScript générés depuis \code{contracts.py}, fixtures régénérées par le vrai moteur \\
    Du code généré par IA faux mais plausible & Tests obligatoires et relecture humaine avant chaque fusion \\
    Calculs trop lents en direct & Monte Carlo vectorisé, cache à chaque étape, budget de calcul fixe \\
    Un YAML importé qui sature le serveur & Plafonds vérifiés au chargement, délais à l'exécution \\
    Un besoin imprévu en pleine nuit & Sorties prévues (section~\ref{sec:imprevu}), sans modifier les contrats \\
    La démo plante & Graine fixée, mode hors réseau, vidéo de secours \\
    L'hébergeur gratuit est lent ou en veille & Mise en ligne dès le premier soir, réveil avant la soutenance, cache, version locale \\
    \og C'est un simulateur de plus \fg{} & Insister sur ce que les autres n'ont pas : QCM, validation par simulation, Correction, plusieurs disciplines, score de Brier \\
    \og Où est l'IA ? \fg{} & Assumer le calcul déterministe (fiable, gratuit, vérifiable) ; l'IA sert à créer des expériences, et le schéma valide tout ce qu'elle produit \\
    \bottomrule
  \end{tabularx}
  \caption{Les principaux risques et leurs garde-fous.}
\end{table}

\textbf{Après le hackathon :} intégration à l'application \emph{Concepts \& Révision}, en réutilisant le moteur déjà en ligne comme bloc Python (les expériences s'insèrent dans les fiches, les QCM rejoignent les révisions espacées), comptes facultatifs pour retrouver son historique sur plusieurs appareils, moteurs \texttt{process} et \texttt{matrix}, et nouvelles expériences issues de la réserve.

\textbf{Hors de la version 1,} et donc hors du hackathon : ces points pourront s'ajouter ensuite sans changer les contrats.
\begin{table}[H]
  \centering
  \small
  \begin{tabularx}{\textwidth}{@{}L{5cm} Y Y@{}}
    \toprule
    \textbf{Point} & \textbf{Pourquoi plus tard} & \textbf{En attendant} \\
    \midrule
    Héritage entre fichiers YAML & résolution et messages d'erreur délicats & ancres YAML (\texttt{\&}, \texttt{<<}) dans un même fichier \\
    Petits multiples (une vue par méthode côte à côte) & champ facultatif, ajouté sans rupture & une vue par méthode si besoin \\
    Mode statique servi par Cloudflare & script de génération, pas un contrat & cache et réveil du serveur avant la soutenance \\
    Calcul dans le navigateur (Pyodide) & gros chantier & serveur Python \\
    Générateur \texttt{process} & déjà en \ptwo{} & fonction de tirage \code{ar1(base, rho)} du registre \\
    Version anglaise & aucune demande pour le jury & textes en français, clés stables \\
    \bottomrule
  \end{tabularx}
  \caption{Ce qui n'entre pas dans la version 1.}
\end{table}

% ============================================================================
% ANNEXES
% ============================================================================
\newpage
\begin{appendices}

\section{Glossaire officiel}\label{ann:glossaire}
Le même mot désigne la même chose dans l'interface, dans les fichiers et dans le code.

\begin{longtable}{@{}L{3.6cm} L{3.6cm} >{\raggedright\arraybackslash}p{\dimexpr\textwidth-7.2cm-4\tabcolsep\relax}@{}}
  \toprule
  \textbf{Terme (interface)} & \textbf{Clé dans le code} & \textbf{Définition} \\
  \midrule
  \endfirsthead
  \toprule
  \textbf{Terme (interface)} & \textbf{Clé dans le code} & \textbf{Définition} \\
  \midrule
  \endhead
  \bottomrule
  \endfoot
  Expérience & \texttt{experiment} & Fichier YAML décrivant un théorème, ses hypothèses, ses explications, ses défis et ses QCM. \\
  Générateur & \texttt{generator} & Processus qui produit les données ou trajectoires du \og vrai monde \fg{}. \\
  Méthode principale & \code{role: method} & Procédure classique du cours (MCO, test de Student, Euler\dots). Une seule par expérience. \\
  Correction & \code{role: correction} & Procédure robuste ou corrigée (HC3, Wilcoxon, Ridge, RK4, Borda). Au plus une. \\
  Cible relative & \code{target: base} & La cible d'une mesure est sa valeur à l'état sain, pour la même méthode (puissance d'un test). \\
  Série d'Entraînement & \texttt{TrainingSession} & 10 questions vérifiées, identifiées par leur filtre et leur graine, partageables par lien. \\
  Hypothèses réparées & \texttt{repairs} & Hypothèses qu'une correction ou une alternative déclare réparer ; vérifié par le test 3. \\
  Alternative & \code{role: alternative} & Autre méthode comparée sur le banc (Newey-West, Huber\dots). Au plus deux. \\
  Comprendre & \texttt{theorem}, \texttt{why}, \code{when_broken} & Panneau qui explique l'énoncé et le rôle de chaque hypothèse. \\
  Mesure & \texttt{measures} & Grandeur calculée par un agrégateur sur les répétitions, pour chaque méthode. \\
  Mesure primaire & \code{primary: true} & La mesure qui pilote le badge, l'équation, la sensibilité et le seuil critique. Une seule par expérience. \\
  Verdict & \texttt{states} & Fiable, Fragile ou Rompu (\texttt{ok}, \texttt{warning}, \texttt{broken}), déduit de la cible et des deux écarts de la mesure. \\
  Erreur Monte Carlo & \code{mc_se} & Incertitude d'une mesure due au nombre fini de tirages ; affichée à $\pm$ 2 erreurs types. \\
  État sain & \code{healthy_value} & Valeur de chaque hypothèse pour laquelle le théorème tient. \\
  Référence & \texttt{baseline} & Résultat de l'état sain, avec les mêmes paramètres et la même graine, superposé en pointillés. \\
  État de calcul & \texttt{EngineState} & Paramètres, hypothèses, points ajoutés et graine : ce qui détermine le résultat. \\
  État d'affichage & \texttt{UiState} & Méthodes affichées, onglet, calques : ne quitte jamais l'interface. \\
  Défis & \texttt{challenges} & Micro-objectifs du Laboratoire, vérifiés par le calcul. \\
  Laboratoire / Entraînement & \texttt{lab} / \texttt{training} & Les deux modes de l'application. \\
  Validation par simulation & \code{reveal_state}, \texttt{verify} & Simulation lancée juste après une réponse QCM, sur l'état qui confirme ou contredit la réponse ; sa condition est vérifiée par le serveur. \\
  Crans de confiance & \texttt{confidence} & Je devine, Plutôt sûr, Certain : traduits en probabilité selon le nombre de choix. \\
  Balayage & \texttt{sweep} & Relance de toutes les méthodes sur une grille d'états, à 1 ou 2 axes. \\
  Étude de sensibilité & \texttt{sensitivity} & Évolution d'une mesure selon une hypothèse (courbe ou barres). \\
  Domaine de validité & \code{validity_domain} & Carte de chaleur croisant deux paramètres. \\
  Répétitions & \texttt{repetitions} & Vue fréquentiste sur 100 à 500 tirages. \\
  Seuil critique & \texttt{critical} & Premier point d'un balayage où la méthode principale quitte le verdict Fiable. \\
  Balayage automatique & \code{auto_sweep} & Animation d'une hypothèse ou de $n$ (fiche~\ref{fiche:balayage}). \\
  Banc de méthodes & \code{methods_on} & Méthodes affichées côte à côte, avec leurs mesures (fiche~\ref{fiche:banc}). \\
  Journal des tirages & \code{draw_log} & Pastilles du test sur le premier échantillon de chaque tirage (fiche~\ref{fiche:journal}). \\
  Inspecteur & \texttt{inspect} & Fiche d'une observation : résidu, levier, distance de Cook (fiche~\ref{fiche:inspecteur}). \\
  Calques & \texttt{layers}, \texttt{layer} & Étapes d'une figure construite couche par couche (fiche~\ref{fiche:calques}). \\
  Contre-exemples & \texttt{counterexamples} & Préréglages chargeant un cas pathologique connu. \\
  Équation dynamique & \texttt{equation} & Formule écrite en gabarit, dont les termes nommés changent de couleur. \\
  Terme d'équation & \code{equation.terms} & Partie nommée de la formule, avec sa version saine et sa version corrigée. \\
  Points ajoutés & \code{extra_points}, \texttt{inject} & Points posés au clic, ajoutés aux variables désignées par le générateur. \\
  Lien de partage & \texttt{share} & URL qui rouvre une expérience dans un état précis. \\
\end{longtable}

\newpage
\section{Anatomie complète d'une expérience (YAML)}\label{ann:yaml}
Trois fichiers de référence servent de modèles, un par expérience P0 : \code{experiments/regression_lineaire.yaml} pour les expériences aléatoires (générateur \texttt{random}), écrit ensemble pendant la première heure ; \code{experiments/test_student.yaml} pour les tests, avec une cible relative ; \code{experiments/pendule_euler.yaml} pour les expériences déterministes (générateur \texttt{ode}). Le premier montre les trois types d'hypothèses, les méthodes en liste, les mesures et leurs verdicts, l'équation à termes nommés, les contre-exemples (dont Anscombe), les défis, les vues et les trois types de questions.

Règles générales : une hypothèse ou un paramètre absent d'un état (contre-exemple, question) prend sa valeur saine ou sa valeur par défaut ; les fonctions \texttt{law}, \texttt{where}, \texttt{sqrt}, \texttt{exp}, \texttt{sin}, \texttt{ar1} et \code{ols_match} ainsi que l'indice d'observation \texttt{idx} font partie de la liste blanche des expressions. Les lignes marquées \texttt{\# P1} peuvent rester dans le fichier dès le départ : elles ne coûtent que leur test.

\subsection{Régression linéaire : le modèle du générateur \texttt{random}}\label{ann:yaml-random}
\begin{lstlisting}[caption={Fichier de référence \texttt{experiments/regression\_lineaire.yaml}.}]
schema_version: 1
id: "regression-lineaire"
title: "Régression linéaire : résidus et validité des intervalles de confiance"
discipline: "statistique"
level: "1A/2A"
concepts: ["mco", "residus", "heteroscedasticite", "levier", "independance", "white-hc3"]

# --- 1. Le théorème, expliqué simplement (panneau « Comprendre ») ---
theorem:
  name: "Gauss-Markov et intervalles de confiance des MCO"
  statement: |
    Sous les hypothèses du modèle linéaire, l'estimateur des MCO est sans biais,
    et l'intervalle de confiance à 95 % contient la vraie pente
    dans 95 % des échantillons.

# --- 2. L'équation dynamique : un gabarit et des termes nommés ---
equation:
  tex: '\widehat{\mathrm{Var}}(\hat{\beta}) = {{var}}'
  terms:
    var:
      ok: '\hat{\sigma}^2 (X^\top X)^{-1}'
      corrected: '(X^\top X)^{-1} X^\top \hat{\Omega}_{\mathrm{HC3}} X (X^\top X)^{-1}'

# --- 3. Les paramètres réglables ---
params:
  n:     { type: number, label: "Taille d'échantillon n", min: 30, max: 1000, step: 10, default: 150,
           sweep_values: [30, 100, 300, 1000] }
  beta1: { type: number, label: "Vraie pente", min: -3.0, max: 3.0, step: 0.1, default: 1.5 }
  match_summary: { type: boolean, label: "Caler la pente et le R2 (Anscombe)", default: false }

# --- 4. Les hypothèses, avec leurs explications ---
hypotheses:
  hetero:
    type: continuous
    range: [0, 1]
    step: 0.05
    healthy_value: 0.0
    label: "Homoscédasticité (variance constante)"
    category: "homoscedasticite"
    ok_label: "Respectée"
    broken_label: "Relâchée"
    statement_tex: '\mathrm{Var}(\varepsilon_i \mid X) = \sigma^2'
    why: "La formule de variance de la pente suppose un même sigma^2 pour toutes les observations."
    when_broken: "La pente reste juste en moyenne, mais les intervalles sont trop étroits : ils couvrent moins de 95 %."
    affects: [var]
  autocorr:                       # P1
    type: boolean
    healthy_value: false
    label: "Indépendance des observations"
    category: "independance"
    ok_label: "Erreurs indépendantes"
    broken_label: "Erreurs et x autocorrélés (AR(1), rho = 0,7)"
    statement_tex: '\mathrm{Cov}(\varepsilon_i, \varepsilon_j) = 0 \ \text{pour } i \neq j'
    why: "La formule de variance suppose que chaque erreur n'apporte aucune information sur ses voisines."
    when_broken: "Si x et les erreurs sont tous deux autocorrélés, les IC classiques et HC3 tombent vers 75 % ; Newey-West remonte vers 87 %."
    affects: [var]
  error_law:
    type: choice
    healthy_value: "normal"
    label: "Loi des erreurs"
    category: "normalite"
    options:
      normal:      { label: "Normale", law: normal }
      exponential: { label: "Exponentielle (asymétrique)", law: exponential }
      student_2:   { label: "Student à 2 ddl (queues lourdes)", law: student_2 }
    statement_tex: '\varepsilon_i \sim \mathcal{N}(0, \sigma^2)'
    why: "La normalité rend les tests de Student exacts quand n est petit."
    when_broken: "Pour la pente des MCO, presque rien : même à n = 30, la couverture reste proche de 95 %. Un résultat à retenir."
    robust_expected: true           # test 2 : la méthode doit rester Fiable
    affects: []
  relation:                       # P1
    type: choice
    healthy_value: "linear"
    label: "Forme de la relation"
    category: "specification"
    options:
      linear:    { label: "Linéaire", expr: "x" }
      quadratic: { label: "Quadratique", expr: "x + 0.5 * (x**2 - 1)" }
      break:     { label: "Rupture de pente", expr: "where(x > 0, 1.8 * x, 0.2 * x)" }
    statement_tex: 'E[Y \mid X] = \beta_0 + \beta_1 X'
    why: "Le modèle suppose que l'espérance de Y varie linéairement avec X."
    when_broken: "La pente reste juste, mais les résidus dessinent une forme et la couverture classique chute (vers 73 % en quadratique)."
    affects: [var]                  # le terme oublié se comporte comme une variance non constante
  add_leverage:
    type: boolean
    healthy_value: false
    label: "Présence d'un point levier"
    category: "outliers"
    ok_label: "Sans levier"
    broken_label: "Avec levier"
    statement_tex: 'h_{ii} \ll 1 \ \text{pour tout } i'
    why: "Aucune observation ne doit peser seule sur l'ajustement."
    when_broken: "Une observation extrême en x attire la droite : la pente change à cause d'un seul point."
    affects: []

# --- 5. Les contre-exemples en un clic ---
counterexamples:
  - id: "grand-n-hetero"
    label: "Grand n mais hétéroscédastique"
    explain: "Même avec n = 600, la couverture des intervalles classiques reste défaillante."
    state: { params: { n: 600 }, hypotheses: { hetero: 0.85 } }
  - id: "erreurs-autocorrelees"
    label: "Erreurs autocorrélées"
    explain: "HC3 ne suffit pas ; l'alternative Newey-West remonte nettement, sans tout réparer à n modéré."
    state: { hypotheses: { autocorr: true } }
  - id: "anscombe-courbe"
    label: "Anscombe : relation courbe"
    explain: "Même pente et même R2 que le cas linéaire, mais des résidus en arc."
    expect: same_summary            # test 4 : pente et R2 identiques, au lieu d'une rupture
    state: { params: { match_summary: true }, hypotheses: { relation: "quadratic" } }
  - id: "anscombe-levier"
    label: "Anscombe : un point levier"
    explain: "Même pente et même R2, mais un seul point porte tout l'ajustement."
    expect: same_summary
    state: { params: { match_summary: true }, hypotheses: { add_leverage: true } }

# --- 6. Le monde simulé ---
generator:
  type: random
  n: p.n
  repeat: 500
  variables:
    x_raw:    { expr: "law('normal')" }
    lev:      { expr: "h.add_leverage & (idx == 0)" }        # la première observation devient le levier
    x_base:   { expr: "where(h.autocorr, ar1(x_raw, 0.7), x_raw)" }
    x:        { expr: "where(lev, 6.0, x_base)" }
    base_eps: { expr: "@h.error_law" }                         # loi de l'option choisie
    eps_ind:  { expr: "base_eps * exp(0.9 * h.hetero * x)" }   # écart type qui croît avec x
    eps:      { expr: "where(h.autocorr, ar1(eps_ind, 0.7), eps_ind)" }
    signal:   { expr: "@h.relation" }                          # pente de projection égale à 1 pour chaque forme
    y_raw:    { expr: "1.0 + p.beta1 * signal + eps + where(lev, -8.0, 0.0)" }
    y:        { expr: "where(p.match_summary, ols_match(x, y_raw, 1.0, p.beta1, 0.67), y_raw)" }
  inject: { x: x, y: y }                                       # variables qui reçoivent les points posés au clic

# --- 7. Les méthodes : principale, correction, alternatives ---
methods:
  - id: ols
    role: method
    call: ols_ci
    label: "MCO classique (IC à 95 %)"
    args: { x: x, y: y, truth: p.beta1, alpha: 0.05, se_type: "classic" }
  - id: hc3
    role: correction
    call: ols_ci
    label: "Écarts-types robustes de White (HC3)"
    explain: "L'estimateur sandwich HC3 estime la variance de chaque résidu séparément."
    repairs: [hetero, relation]
    args: { x: x, y: y, truth: p.beta1, alpha: 0.05, se_type: "HC3" }
  - id: nw                          # P1 (banc de méthodes)
    role: alternative
    call: ols_ci
    label: "Écarts-types de Newey-West (HAC)"
    explain: "L'estimateur HAC tient compte de la variance non constante et de la corrélation entre erreurs voisines."
    repairs: [hetero, autocorr]
    args: { x: x, y: y, truth: p.beta1, alpha: 0.05, se_type: "HAC" }

# --- 8. Les mesures et leurs verdicts ---
measures:
  coverage:
    agg: "rate(reps.hit)"
    label: "Taux de couverture réel"
    unit: "%"
    scale: 100
    decimals: 1
    primary: true
    target: 0.95
    direction: below            # below, above ou both : le côté où l'écart est grave
    ok_within: 0.03             # Fiable à partir de 92 %
    broken_beyond: 0.10         # Rompu sous 85 %
    texts:
      ok: "Intervalle fiable : {value} de couverture."
      warning: "Couverture en baisse ({value}) : la variance commence à être sous-estimée."
      broken: "Intervalle non fiable ({value}) : la variance est sous-estimée."
  bias:
    agg: "bias(reps.est, p.beta1)"
    label: "Biais de la pente"
    decimals: 3
    target: 0.0
    direction: both
    ok_within: 0.05
    broken_beyond: 0.15
    texts:
      ok: "Pente sans biais ({value})."
      warning: "Pente légèrement biaisée ({value})."
      broken: "Pente biaisée ({value})."
  width:
    agg: "mean(reps.width)"
    label: "Largeur moyenne des IC"    # sans cible : sert à lire le coût d'une méthode robuste
    decimals: 3
  bp_pvalue:
    agg: "first(reps.bp_pvalue)"
    label: "p-value de Breusch-Pagan (échantillon affiché)"
    decimals: 3                 # sans cible : mesure informative, sans verdict

# --- 9. Les défis ---
challenges:
  - id: "trap-bp-test"
    goal: "Piégez le test de Breusch-Pagan : sur l'échantillon affiché, obtenez une p-value > 0,05 alors que la couverture est rompue."
    when: "m.bp_pvalue > 0.05 and s.coverage == 'broken'"
    hint: "Essayez un petit n et plusieurs tirages : à n = 30 et h = 0,7, environ un échantillon sur cinq trompe le test."
    unlock: "Vous illustrez l'angle mort d'un test peu puissant : ne pas rejeter H0 ne prouve pas que l'hypothèse tient."
  - id: "big-n-limit"
    goal: "Vérifiez qu'augmenter n au-delà de 500 ne répare pas les IC classiques quand hetero >= 0,7."
    when: "p.n >= 500 and h.hetero >= 0.7 and s.coverage == 'broken'"
    unlock: "L'erreur de variance des MCO sous hétéroscédasticité ne disparaît pas quand n grandit."

# --- 10. Les vues et les mécaniques visuelles ---
auto_sweep: [h.hetero, p.n]       # boutons Lecture du balayage automatique (fiche F.3)

views:
  - id: nuage
    type: scatter
    title: "Nuage et droite ajustée"
    allow_point_injection: true     # fiche F.4
    inspect: true                   # fiche F.6
    axes: { x: auto, y: auto }
    layers:                         # fiche F.7, dans l'ordre des étapes
      - { id: nuage, label: "Le nuage" }
      - { id: sct,   label: "Écarts à la moyenne (SCT)" }
      - { id: sce,   label: "Écarts expliqués (SCE)" }
      - { id: scr,   label: "Résidus (SCR)" }
    bind:
      x: { path: example.x, layer: nuage }
      y: { path: example.y, layer: nuage }
      fit: { path: example.y_hat, layer: sce }
      segments:
        - { from: example.y,     to: example.y_mean, layer: sct, label: "y - moyenne" }
        - { from: example.y_hat, to: example.y_mean, layer: sce, label: "prédiction - moyenne" }
        - { from: example.y,     to: example.y_hat,  layer: scr, label: "résidu" }
  - id: residus
    type: residuals
    title: "Résidus"
    bind: { x: example.x, resid: example.resid }
  - id: ic
    type: repetitions
    mode: intervals
    title: "100 intervalles de confiance empilés"
    bind: { est: reps.est, low: reps.low, high: reps.high, hit: reps.hit, truth: p.beta1 }
  - id: sens
    type: sensitivity
    title: "Étude de sensibilité"
    sweep: { axes: [h.hetero] }
  - id: validite
    type: validity_domain
    title: "Domaine de validité"
    sweep: { axes: [h.hetero, p.n] }

# --- 11. Les questions d'Entraînement ---
opening_question: "reg-q1-biais"   # P1 : carte « Pariez » non bloquante dans le Laboratoire

questions:
  - id: "reg-q1-biais"
    type: misconception           # 2 choix : crans à 50 %, 72,5 % et 95 %
    prompt: "Sous forte hétéroscédasticité (n = {n}), que devient l'estimateur MCO de la pente ?"
    hypothesis: "hetero"          # explication affichée après la réponse
    procedural_params:
      n: { min: 150, max: 600, step: 50 }
    choices:
      - id: "a"
        label: "Il devient biaisé, même pour un grand échantillon."
        correct: false
        reveal_state: { params: { n: "{n}" }, hypotheses: { hetero: 0.9 } }
        verify: "s.bias == 'ok' and s.coverage == 'broken'"
        focus_view: "ic"
        explain: "Observez les centres des 100 intervalles : ils restent centrés sur la vraie pente. Seule la largeur est fausse."
      - id: "b"
        label: "Il reste sans biais, mais les intervalles de confiance sont invalides."
        correct: true
        reveal_state: { params: { n: "{n}" }, hypotheses: { hetero: 0.9 } }
        verify: "s.bias == 'ok' and s.coverage == 'broken'"
        focus_view: "ic"
        explain: "Exact : l'estimation reste juste en moyenne tandis que la couverture chute."

  - id: "reg-q2-diag"
    type: diagnostic_auto         # 4 choix : crans à 25 %, 60 % et 95 %
    prompt: "Analysez le graphique des résidus ci-dessous. Quelle situation observez-vous ?"
    view: "residus"
    pathologies:
      from: [hetero, add_leverage, relation.quadratic]   # plus « Hypothèses respectées », ajouté d'office

  - id: "reg-q3-n"                # P1
    type: prediction
    prompt: "Sous forte hétéroscédasticité, que devient la couverture des IC classiques si n passe de {n} à {n} x 2 ?"
    procedural_params:
      n: { min: 100, max: 400, step: 50 }
    state: { hypotheses: { hetero: 0.9 } }
    sweep: { target: p.n, from: "{n}", to: "{n} * 2" }
    measure: coverage             # choix fixes : augmente, diminue, reste stable ; réponse calculée par le moteur
\end{lstlisting}

\subsection{Test de Student : une cible relative et deux hypothèses \texorpdfstring{$H_0$ et $H_1$}{H0 et H1}}\label{ann:yaml-student}
Ce fichier montre ce qui change pour un test : le générateur simule le même échantillon sous $H_0$ et sous $H_1$ (décalé de $\delta$), la méthode renvoie les rejets sous les deux, et la puissance est jugée par rapport à la loi normale grâce à \code{target: base}. Les chiffres des commentaires viennent d'une simulation faite avant le hackathon (500 répétitions).

\begin{lstlisting}[caption={Fichier de référence \texttt{experiments/test\_student.yaml}.}]
schema_version: 1
id: "test-student"
title: "Test de Student : queues lourdes, niveau et puissance"
discipline: "statistique"
level: "1A"
concepts: ["test-de-student", "wilcoxon", "queues-lourdes", "puissance", "risque-alpha"]

theorem:
  name: "Test de Student pour un échantillon"
  statement: |
    Si les observations sont indépendantes et de loi normale, la statistique
    T = moyenne / (S / racine de n) suit une loi de Student à n - 1 degrés de liberté :
    le test rejette H0 à tort dans exactement 5 % des cas.

equation:
  tex: '{{stat}}'
  terms:
    stat:
      ok: 'T = \frac{\bar{X}}{S / \sqrt{n}} \sim t_{n-1}'
      corrected: 'W^{+} = \sum_{i} R_i \, \mathbf{1}\{X_i > 0\}'

params:
  n:     { type: number, label: "Taille d'échantillon n", min: 5, max: 200, step: 1, default: 25,
           sweep_values: [5, 10, 25, 50, 100, 200] }
  delta: { type: number, label: "Décalage sous H1", min: 0.1, max: 1.5, step: 0.1, default: 0.5 }

hypotheses:
  loi:
    type: choice
    healthy_value: "normal"
    label: "Loi des observations (toutes symétriques)"
    category: "normalite"
    options:
      normal:    { label: "Normale", law: normal }
      uniform:   { label: "Uniforme", law: uniform, robust_expected: true }
      bimodal:   { label: "Bimodale", law: bimodal, robust_expected: true }
      student_2: { label: "Student à 2 ddl (queues lourdes)", law: student_2 }
      cauchy:    { label: "Cauchy (sans espérance)", law: cauchy }
    statement_tex: 'X_i \sim \mathcal{N}(\mu, \sigma^2)'
    why: "La normalité donne la loi exacte de T ; la moyenne et l'écart type empiriques supposent des queues légères."
    when_broken: "Avec des queues lourdes, un seul point extrême gonfle S : le test devient trop prudent et perd presque toute sa puissance."
    affects: [stat]

counterexamples:
  - id: "cauchy-petit-n"
    label: "Cauchy, n = 25"
    explain: "Le test de Student ne détecte presque plus rien : comparez sa puissance à celle de la Référence (loi normale)."
    state: { hypotheses: { loi: "cauchy" } }
  - id: "cauchy-grand-n"
    label: "Cauchy, n = 200"
    explain: "Un grand échantillon ne sauve pas le test de Student ; Wilcoxon retrouve toute sa puissance."
    state: { params: { n: 200 }, hypotheses: { loi: "cauchy" } }

generator:
  type: random
  n: p.n
  repeat: 500
  variables:
    x0: { expr: "@h.loi" }               # sous H0 : centre de symétrie en 0
    x1: { expr: "x0 + p.delta" }         # sous H1 : mêmes tirages, décalés

methods:
  - id: student
    role: method
    call: one_sample_test
    label: "Test de Student"
    args: { x0: x0, x1: x1, alpha: 0.05, kind: "student" }
  - id: wilcoxon
    role: correction
    call: one_sample_test
    label: "Test de Wilcoxon (rangs signés)"
    explain: "Les rangs ne voient pas la taille des valeurs extrêmes : sous une loi symétrique, le niveau reste exact et la puissance tient."
    repairs: [loi]
    args: { x0: x0, x1: x1, alpha: 0.05, kind: "wilcoxon" }

measures:
  power:
    agg: "rate(reps.reject_h1)"
    label: "Puissance sous H1"
    unit: "%"
    scale: 100
    decimals: 1
    primary: true
    target: base                 # la cible est la puissance de la même méthode sous loi normale
    direction: below
    ok_within: 0.10
    broken_beyond: 0.25          # n = 25 : Student Rompu sous Student à 2 ddl et Cauchy ; Wilcoxon Fragile sous Cauchy
    texts:
      ok: "Puissance préservée : {value}."
      warning: "Puissance entamée : {value}."
      broken: "Puissance effondrée : {value}."
  alpha:
    agg: "rate(reps.reject)"
    label: "Risque de première espèce réel"
    unit: "%"
    scale: 100
    decimals: 1
    target: 0.05
    direction: both
    ok_within: 0.025             # au moins 2,5 erreurs Monte Carlo : le hasard seul ne change pas le verdict
    broken_beyond: 0.04          # Cauchy : 1,8 % à n = 25 (Fragile), 0,4 % à n = 200 (Rompu)

draw_log: { reject: reps.reject, pvalue: reps.pvalue, reject_label: "faux positif" }   # P1

challenges:
  - id: "moitie-puissance"
    goal: "Trouvez une loi pour laquelle le test de Student perd plus de la moitié de sa puissance."
    when: "m.power < 0.5 * base.power"
    unlock: "Avec des queues lourdes, l'écart type empirique explose et masque l'effet."
  - id: "grand-n-cauchy"
    goal: "Montrez qu'un grand échantillon ne rend pas sa puissance au test de Student sous Cauchy."
    when: "h.loi == 'cauchy' and p.n >= 100 and s.power == 'broken'"
    unlock: "Sans espérance, il n'y a pas de théorème central limite : n ne répare rien."

auto_sweep: [p.n]

views:
  - id: pvalues
    type: distribution
    title: "Distribution des p-values sous H0"
    bind: { stat: reps.pvalue }
    ui: { x_range: [0, 1], uniform_reference: true }   # un test exact donne des p-values uniformes
  - id: echantillon
    type: distribution
    title: "Un échantillon"
    bind: { stat: example.x0 }
  - id: sens
    type: sensitivity
    title: "Puissance selon la loi"
    sweep: { axes: [h.loi] }
  - id: validite
    type: validity_domain
    title: "Loi et taille d'échantillon"
    sweep: { axes: [h.loi, p.n] }

questions:
  - id: "stu-q1-cauchy"
    type: misconception
    prompt: "Sous une loi de Cauchy (n = {n}), le test de Student rejette-t-il H0 à tort plus souvent que 5 % ?"
    hypothesis: "loi"
    procedural_params:
      n: { min: 15, max: 60, step: 5 }
    choices:
      - id: "a"
        label: "Oui, beaucoup plus souvent."
        correct: false
        reveal_state: { params: { n: "{n}" }, hypotheses: { loi: "cauchy" } }
        verify: "m.alpha < 0.05 and s.power == 'broken'"
        focus_view: "pvalues"
        explain: "Les p-values s'entassent près de 1 : le test devient trop prudent, et c'est sa puissance qui s'effondre."
      - id: "b"
        label: "Non : il devient trop prudent, et c'est sa puissance qui s'effondre."
        correct: true
        reveal_state: { params: { n: "{n}" }, hypotheses: { loi: "cauchy" } }
        verify: "m.alpha < 0.05 and s.power == 'broken'"
        focus_view: "pvalues"
        explain: "Exact : un point extrême gonfle S, donc T reste petit, sous H0 comme sous H1."

  - id: "stu-q2-diag"
    type: diagnostic_auto
    prompt: "Voici un échantillon de 200 observations. De quelle loi vient-il ?"
    view: "echantillon"
    state: { params: { n: 200 } }
    pathologies:
      from: [loi.student_2, loi.cauchy]      # plus « Normale », ajoutée d'office
\end{lstlisting}

La méthode \code{one_sample_test} déclare \code{provides={"reps": ["reject", "reject_h1", "pvalue", "stat"], "example": ["x0"]}}. Elle est vectorisée : \code{scipy.stats.ttest_1samp} et \code{scipy.stats.wilcoxon} acceptent l'argument \texttt{axis} et traitent les 500 échantillons d'un coup.

\subsection{Pendule et méthode d'Euler : le modèle du générateur \texttt{ode}}\label{ann:yaml-ode}
Une expérience déterministe n'a pas de répétitions : le générateur \texttt{ode} calcule une seule fois la solution de référence et fournit le système aux méthodes, qui l'intègrent avec leur propre schéma. Les mesures passent par \code{summary(k)} ; elles n'ont pas d'erreur Monte Carlo, et l'interface n'affiche pas de marge.

Ce que fournit le générateur \texttt{ode}, dans l'espace de noms des expressions et des arguments :
\begin{itemize}[leftmargin=*, nosep]
  \item \texttt{t} : la grille d'affichage (400 points de 0 à \code{t_end}) ;
  \item une variable par état (\texttt{theta}, \texttt{omega}) : la solution de référence sur cette grille, obtenue par \code{solve_ivp} avec une tolérance de $10^{-10}$ ;
  \item \texttt{system} : l'objet qui porte le second membre, les conditions initiales et \code{t_end}, passé tel quel aux méthodes ;
  \item les variables \texttt{derived}, calculées une fois (ici le pas \texttt{dt}).
\end{itemize}
Une méthode d'intégration renvoie sa trajectoire interpolée sur la grille \texttt{t} (dans \texttt{curves}) et ses grandeurs globales (dans \texttt{summary}). Les chiffres des commentaires viennent d'une simulation faite avant le hackathon ($\theta_0 = 1{,}2$ rad, 5 secondes).

\begin{lstlisting}[caption={Fichier de référence \texttt{experiments/pendule\_euler.yaml}.}]
schema_version: 1
id: "pendule-euler"
title: "Pendule simple : la méthode d'Euler et son pas de temps"
discipline: "physique"
level: "1A"
concepts: ["edo", "euler", "runge-kutta", "stabilite-numerique", "pendule"]

theorem:
  name: "Convergence de la méthode d'Euler explicite"
  statement: |
    Sur un intervalle de temps fixé, l'erreur de la méthode d'Euler est proportionnelle
    au pas de temps : elle tend vers 0 quand le pas tend vers 0, mais elle s'accumule
    vite quand le pas grandit.

equation:
  tex: '\max_k |\theta(t_k) - \theta_k| = {{err}}'
  terms:
    err: { ok: 'O(\Delta t)', corrected: 'O(\Delta t^4)' }

params:
  theta0: { type: number, label: "Amplitude initiale (rad)", min: 0.2, max: 2.5, step: 0.1, default: 1.2,
            sweep_values: [0.2, 0.6, 1.2, 2.0] }
  t_end:  { type: number, label: "Durée simulée (s)", min: 2, max: 10, step: 1, default: 5 }

hypotheses:
  pas:
    type: continuous
    healthy_value: 0.0
    label: "Pas de temps petit"
    category: "stabilite-numerique"
    ok_label: "Pas fin"
    broken_label: "Pas grossier"
    statement_tex: '\Delta t \ll T_0'
    why: "Euler suit la tangente pendant tout un pas : l'erreur commise à chaque pas s'ajoute à la précédente."
    when_broken: "L'énergie du pendule augmente à chaque pas : les oscillations s'amplifient, jusqu'à faire faire le tour au pendule."
    affects: [err]

counterexamples:
  - id: "tour-complet"
    label: "Le pendule fait le tour"
    explain: "Avec un pas de 0,06 s, Euler fait passer le pendule par-dessus son axe ; RK4 reste exact."
    state: { hypotheses: { pas: 0.9 } }

generator:
  type: ode
  t_end: p.t_end
  state: { theta: p.theta0, omega: 0.0 }      # variables d'état et conditions initiales
  rhs:
    theta: "omega"
    omega: "-9.81 * sin(theta)"               # g / L avec L = 1 m
  derived:
    dt: "10 ** (-3.5 + 2.5 * h.pas)"          # de 0,0003 s (h = 0) à 0,1 s (h = 1), échelle logarithmique

methods:
  - id: euler
    role: method
    call: integrate
    label: "Euler explicite"
    args: { system: system, dt: dt, scheme: "euler" }
  - id: rk4
    role: correction
    call: integrate
    label: "Runge-Kutta d'ordre 4"
    explain: "RK4 évalue la pente quatre fois par pas : l'erreur devient proportionnelle à la puissance 4 du pas."
    repairs: [pas]
    args: { system: system, dt: dt, scheme: "rk4" }

measures:
  err_angle:
    agg: "summary(max_rel_error)"     # max |theta - theta_ref| / theta0
    label: "Erreur maximale sur l'angle"
    unit: "%"
    scale: 100
    decimals: 1
    primary: true
    target: 0.0
    direction: above
    ok_within: 0.05                   # Euler Fiable jusqu'à h = 0,2 environ (1,2 % à h = 0)
    broken_beyond: 0.20               # Rompu à partir de h = 0,5 environ (45 % à h = 0,6)
    texts:
      ok: "Trajectoire fidèle : erreur de {value}."
      warning: "La trajectoire dérive : erreur de {value}."
      broken: "Trajectoire fausse : erreur de {value}."
  energy_drift:
    agg: "summary(energy_drift)"
    label: "Dérive de l'énergie"
    unit: "%"
    scale: 100
    decimals: 1
  full_turn:
    agg: "summary(full_turn)"         # 1 si le pendule dépasse pi, 0 sinon (Euler : à partir de h = 0,9)
    label: "Le pendule fait le tour"
    decimals: 0

challenges:
  - id: "plus-grand-pas"
    goal: "Trouvez un pas au moins trois fois plus grand que le pas minimal pour lequel Euler reste fiable."
    when: "h.pas >= 0.2 and s.err_angle == 'ok'"
    unlock: "Le pas acceptable dépend de la durée simulée : doublez la durée et recommencez."
  - id: "rk4-gros-pas"
    goal: "Montrez que RK4 reste fiable avec un pas pour lequel Euler fait faire le tour au pendule."
    when: "m.full_turn == 1 and m.err_angle.rk4 < 0.05"
    unlock: "Un schéma d'ordre 4 autorise des pas bien plus grands pour la même précision."

auto_sweep: [h.pas]

views:
  - id: trajectoire
    type: curves
    title: "Angle au cours du temps"
    bind: { t: curves.t, series: [ { y: curves.theta, label: "Angle" } ], truth: gen.theta }
  - id: sens
    type: sensitivity
    title: "Erreur selon le pas"
    sweep: { axes: [h.pas] }
  - id: validite
    type: validity_domain
    title: "Pas et amplitude"
    sweep: { axes: [h.pas, p.theta0] }

questions:
  - id: "pend-q1-rk4"
    type: misconception
    prompt: "RK4 a-t-il besoin d'un pas aussi fin qu'Euler pour rester fiable ?"
    hypothesis: "pas"
    choices:
      - id: "a"
        label: "Oui : toutes les méthodes ont besoin du même pas."
        correct: false
        reveal_state: { hypotheses: { pas: 0.8 } }
        verify: "s.err_angle == 'broken' and m.err_angle.rk4 < 0.05"
        focus_view: "trajectoire"
        explain: "Avec ce pas, Euler s'effondre alors que RK4 suit encore la solution de référence."
      - id: "b"
        label: "Non : RK4 reste fiable avec un pas bien plus grand."
        correct: true
        reveal_state: { hypotheses: { pas: 0.8 } }
        verify: "s.err_angle == 'broken' and m.err_angle.rk4 < 0.05"
        focus_view: "trajectoire"
        explain: "Exact : l'erreur de RK4 est proportionnelle à la puissance 4 du pas."
\end{lstlisting}

La méthode \texttt{integrate} déclare \code{provides={"curves": ["t", "theta"], "summary": ["max_rel_error", "energy_drift", "full_turn"]}} et \code{cost="t_end / dt"}. Une intégration de 16\,000 pas avec Euler prend environ 12 ms en Python : le balayage complet (11 pas, 2 méthodes) tient en moins d'une demi-seconde.

\subsection{Le générateur \texttt{function} \ptwo{}}\label{ann:yaml-function}
Son format est fixé dès maintenant, pour que l'expérience TFA puisse s'écrire plus tard sans toucher au schéma. Il évalue des fonctions sur une grille, sans répétitions ; les mesures passent par \code{summary(k)}, comme pour \texttt{ode}.

\begin{lstlisting}[caption={Extrait du générateur \texttt{function}.}]
generator:
  type: function
  grid: { x: [-3.0, 3.0, 400] }       # minimum, maximum, nombre de points
  variables:
    f: { expr: "@h.forme" }           # options de l'hypothèse : expressions en x
\end{lstlisting}


\newpage
\section{Contrats d'API}\label{ann:api}

\subsection{Les routes}
\begin{table}[H]
  \centering
  \footnotesize
  \begin{tabularx}{\textwidth}{@{}l l c l Y@{}}
    \toprule
    \textbf{Verbe} & \textbf{Route} & \textbf{Prio.} & \textbf{Renvoie} & \textbf{Rôle} \\
    \midrule
    \texttt{GET}  & \code{/api/experiments}        & P0 & \code{ExperimentSummary[]} & Liste des expériences officielles, filtrable. \\
    \texttt{GET}  & \code{/api/experiments/{id}}   & P0 & \texttt{ExperimentSpec}    & Contrôles, explications, équation, défis, contre-exemples. \\
    \texttt{POST} & \code{/api/run}                & P0 & \texttt{RunResult}         & Exécute une expérience officielle (\code{experiment_id}) ou partagée (\code{raw_content}). \\
    \texttt{POST} & \code{/api/sweep}              & P0 / P1 & \texttt{SweepResult}  & Balayage à 1 axe (P0 : sensibilité, $n \to \infty$) ou à 2 axes (P1 : Domaine de validité). \\
    \texttt{GET}  & \code{/api/training/next}      & P0 & \code{TrainingSession} & Prend une série prête dans la réserve d'un filtre et lance son remplacement en arrière-plan (annexe~\ref{ann:api-training}). \\
    \texttt{GET}  & \code{/api/training/session}   & P0 & \code{TrainingSession} & Rouvre une série précise (lien partagé), refabriquée à l'identique si besoin. \\
    \texttt{POST} & \code{/api/validate}           & P1 & \texttt{ValidateResponse}  & Valide un YAML, lance les tests scientifiques et renvoie les erreurs par numéro de ligne. \\
    \bottomrule
  \end{tabularx}
  \caption{Les 7 routes du serveur FastAPI. L'interface les appelle à l'adresse \texttt{VITE\_API\_URL} ; le partage passe directement par Supabase.}
\end{table}

\subsection{Le partage par lien}
\begin{itemize}[leftmargin=*, nosep]
  \item \textbf{Expérience officielle :} l'interface encode l'état de calcul et l'état d'affichage dans l'URL (\code{?exp=...&v=1&p.n=...&h.hetero=...&seed=...&on=...&pts=...}) et les relit à l'ouverture. Aucun serveur n'est impliqué.
  \item \textbf{Expérience importée :} une table Supabase \code{shared_experiments(id, schema_version, content, content_hash, created_at)} ; insertion anonyme, lecture publique par identifiant, aucune modification après insertion. L'interface récupère le YAML puis appelle \code{/api/run} avec \code{raw_content}.
\end{itemize}

\subsection{Le contrat principal : \texttt{POST /api/run}}
Le moteur calcule toujours toutes les méthodes et la Référence : la requête ne dit pas quoi afficher. Toutes les valeurs sont données par méthode.

\begin{lstlisting}[caption={Requête et réponse de \texttt{POST /api/run}.}]
// Requete (RunRequest) : experiment_id OU raw_content (YAML d'une experience partagee)
{
  "experiment_id": "regression-lineaire",
  "raw_content": null,
  "state": {
    "params": { "n": 150, "beta1": 1.5 },
    "hypotheses": { "hetero": 0.8, "error_law": "normal", "add_leverage": false },
    "extra_points": [{ "x": 4.2, "y": -2.0 }],
    "seed": 42
  }
}

// Reponse (RunResult)
{
  "schema_version": 1,
  "state_hash": "a3f09c1e7b2d4e88",
  "methods": [
    { "id": "ols", "role": "method",      "label": "MCO classique (IC à 95 %)" },
    { "id": "hc3", "role": "correction",  "label": "HC3", "explain": "..." },
    { "id": "nw",  "role": "alternative", "label": "Newey-West", "explain": "..." }
  ],
  "measures": {
    "coverage": {
      "label": "Taux de couverture réel", "unit": "%", "scale": 100, "decimals": 1,
      "primary": true, "target": 0.95, "direction": "below",
      "ok_within": 0.03, "broken_beyond": 0.10,
      "values":   { "ols": 0.740, "hc3": 0.946, "nw": 0.941 },
      "mc_se":    { "ols": 0.020, "hc3": 0.010, "nw": 0.011 },
      "states":   { "ols": "broken", "hc3": "ok", "nw": "ok" },
      "texts":    { "ols": "Intervalle non fiable (74,0 %) : ...", "hc3": "...", "nw": "..." },
      "baseline": { "ols": 0.952, "hc3": 0.950, "nw": 0.948 }
    }
  },
  "equation": { "terms": { "var": { "violated_by": ["hetero"] } } },
  "challenges": [ { "id": "trap-bp-test", "passed": false } ],
  "views": [
    { "id": "nuage", "type": "scatter", "title": "Nuage et droite ajustée",
      "axes": { "x": [-3.1, 6.4], "y": [-9.8, 11.4] }, "shown": 151, "total": 151, "clipped": 0,
      "data": { "...": "voir annexe E" } }
  ],
  "draw_log": null,              // test du premier échantillon si l'expérience déclare draw_log
  "invalid_rate": { "ols": 0.0, "hc3": 0.0, "nw": 0.0 },
  "warnings": [ { "code": "repeat_reduced", "text": "Estimation sur 200 tirages pour tenir le délai." } ],
  "repeat_used": 500,
  "compute_ms": 26.4
}
\end{lstlisting}

L'équation renvoie seulement, pour chaque terme, les hypothèses cassées qui le mettent en cause (\code{violated_by}). L'interface applique la règle de couleur de la section~\ref{sec:verdict} avec les méthodes affichées, sans nouvel appel.

\subsection{Le balayage : \texttt{POST /api/sweep}}
Un axe porte sur une hypothèse (\code{h.<id>}) ou sur un paramètre (\code{p.<id>}). Valeurs par défaut : 11 points de 0 à 1 pour une hypothèse continue, toutes les options pour un choix, les deux valeurs pour un interrupteur, \code{sweep_values} pour un paramètre (un paramètre sans ce champ n'est pas balayable). Le reste de l'état ne change pas, et la même graine sert dans toutes les cases. La mesure balayée est la mesure primaire par défaut. Avec 2 axes, \texttt{values}, \code{mc_se} et \texttt{states} deviennent des tableaux à deux dimensions, indexés par l'axe 2 puis l'axe 1.

\begin{lstlisting}[caption={Requête et réponse de \texttt{POST /api/sweep}.}]
// Requete (SweepRequest)
{
  "experiment_id": "regression-lineaire",
  "state": { "params": { "n": 150 }, "hypotheses": { "hetero": 0.8 }, "seed": 42 },
  "axes": [ { "target": "h.hetero" } ],
  "measure": "coverage"
}

// Reponse (SweepResult)
{
  "axes": [ { "target": "h.hetero", "values": [0, 0.1, 0.2, "..."], "labels": ["0", "0,1", "0,2", "..."] } ],
  "measure": "coverage",
  "values": { "ols": [0.951, 0.948, 0.937, "..."], "hc3": [0.949, 0.950, 0.948, "..."] },
  "mc_se":  { "ols": [0.010, 0.010, 0.011, "..."], "hc3": [0.010, 0.010, 0.010, "..."] },
  "states": { "ols": ["ok", "ok", "ok", "..."],    "hc3": ["ok", "ok", "ok", "..."] },
  "critical": { "ols": { "warning_at": 0.27, "broken_at": 0.58 }, "hc3": null },
  "target": 0.95, "ok_within": 0.03, "broken_beyond": 0.10, "direction": "below",
  "current": { "h.hetero": 0.8 },
  "repeat_used": 500, "reduced": false, "compute_ms": 2310
}
\end{lstlisting}

\texttt{critical} n'existe que pour un seul axe continu ; il est interpolé linéairement entre deux points de la grille.

\subsection{L'Entraînement : les séries}\label{ann:api-training}
Une série est fabriquée d'un coup et vérifiée question par question (section~\ref{sec:fabrique}). Le serveur renvoie tout ce qu'il faut pour la jouer, réponses comprises : c'est un outil de révision, pas un examen, et cela évite un aller-retour par question.

Le \emph{filtre} vaut \texttt{toutes}, un identifiant d'expérience (\texttt{regression-lineaire}) ou \code{discipline:<nom>}. Une série de \texttt{toutes} compte 10 questions ; une série d'une seule expérience reprend ses 3 à 5 modèles de questions avec d'autres paramètres tirés au hasard, jusqu'à 10 questions.

\begin{lstlisting}[caption={Les deux routes de l'Entraînement et leur réponse commune.}]
// Serie prete, prise dans la reserve (reponse immediate)
GET /api/training/next?filter=regression-lineaire

// Serie precise, ouverte depuis un lien partage
GET /api/training/session?filter=regression-lineaire&seed=123&h=3f09c1

// Reponse (TrainingSession)
{
  "id": "regression-lineaire~123",
  "filter": "regression-lineaire",
  "seed": 123,
  "content_hash": "3f09c1",
  "stale_link": false,              // true si h ne correspond plus au contenu actuel des YAML
  "questions": [
    {
      "qid": "regression-lineaire/reg-q2-diag/0",
      "experiment_id": "regression-lineaire",
      "type": "diagnostic_auto",
      "prompt": "Analysez le graphique des résidus ci-dessous. Quelle situation observez-vous ?",
      "initial_view": { "id": "residus", "type": "residuals", "...": "format de l'annexe E" },
      "hypothesis": "hetero",
      "k": 4,
      "confidence": { "guess": 0.25, "fair": 0.60, "sure": 0.95 },
      "choices": [
        { "id": "c1", "label": "Hétéroscédasticité", "correct": true,
          "explain": "Vous observez une dispersion en entonnoir caractéristique.",
          "reveal": { "state": { "...": "état de calcul complet" }, "focus_view": "ic" } },
        { "id": "c2", "label": "Point levier influent", "correct": false, "explain": "...", "reveal": { "...": "..." } }
      ]
    }
  ],
  "dropped": [ { "experiment_id": "...", "question_id": "...", "reason": "verify a échoué 5 fois" } ]
}
\end{lstlisting}

L'interface écrit \texttt{train}, \texttt{seed}, \texttt{v} et \texttt{h} dans l'URL dès qu'une série est ouverte : le lien affiché est toujours partageable. L'historique du Bilan est rangé par identifiant de série dans \texttt{localStorage}.

\textbf{Déroulé d'une question côté interface :}
\begin{enumerate}[leftmargin=*, nosep]
  \item afficher l'énoncé et \code{initial_view} (ou rien pour une question en texte), puis les choix dans l'ordre reçu, déjà mélangé ;
  \item l'étudiant choisit une réponse et un cran ; la probabilité vient de \texttt{confidence}, et la réponse est enregistrée dans \texttt{localStorage} pour le Bilan ;
  \item l'interface appelle \code{/api/run} avec \code{reveal.state} du choix sélectionné, ou du bon choix si la réponse est fausse. Le serveur a déjà calculé cet état pendant la vérification : la réponse sort du cache ;
  \item elle affiche la vue \code{focus_view}, \texttt{explain} et le \code{when_broken} de l'hypothèse ;
  \item \og Ouvrir dans le Laboratoire \fg{} ouvre l'URL de \code{reveal.state}, avec le panneau Comprendre sur \texttt{hypothesis}.
\end{enumerate}
Pour une question de prédiction \pone{}, les choix sont toujours \og augmente \fg{}, \og diminue \fg{} et \og reste stable \fg{}, et \texttt{correct} est calculé par le serveur à partir du balayage.

\subsection{La validation d'un YAML : \texttt{POST /api/validate} \pone{}}
\begin{lstlisting}[caption={Réponse de \texttt{POST /api/validate} (la requête contient seulement \texttt{raw\_content}).}]
{
  "ok": false,
  "errors": [
    { "line": 42, "path": "hypotheses.hetero.affects",
      "message": "Le terme « var2 » n'existe pas dans l'équation." }
  ],
  "warnings": [
    { "line": 88, "path": "measures.coverage.ok_within",
      "message": "La tolérance est inférieure à deux erreurs Monte Carlo : le verdict sera instable." }
  ],
  "tests": [
    { "id": 1, "name": "Le théorème tient à l'état sain", "passed": true, "detail": null },
    { "id": 3, "name": "Chaque correction répare ce qu'elle annonce", "passed": false,
      "detail": "hetero = 1 : HC3 reste Rompu (couverture 81,2 %)." }
  ]
}
\end{lstlisting}
Les tests ne sont lancés que si le schéma est valide, avec 200 répétitions pour répondre en quelques secondes.

\newpage
\section{Contrat entre le moteur et les méthodes}\label{ann:methodes}
R2 écrit les méthodes le jour, R1 les branche dans le moteur la nuit. Ce contrat, écrit dans \code{backend/engine/method_contract.py} pendant les 2 premières heures, garantit que les deux parties s'emboîtent sans adaptation.

\subsection{Ce qu'une méthode reçoit et déclare}
Une méthode est une fonction Python pure, enregistrée par un décorateur. Le moteur résout d'abord les arguments écrits dans le YAML (\texttt{x: x} devient le tableau \texttt{x} du générateur, \texttt{truth: p.beta1} devient la valeur du paramètre), puis appelle la fonction avec ces arguments nommés. Le décorateur déclare aussi ce que la méthode fournit et son coût :

\begin{lstlisting}[language=Python, caption={Signature type d'une méthode.}]
@register_method(
    "ols_ci",
    provides={
        "reps": ["est", "low", "high", "hit", "width", "bp_pvalue"],
        "example": ["x", "y", "y_hat", "y_mean", "resid", "leverage", "cooks_d"],
    },
    cost="repeat * n",
)
def ols_ci(*, x, y, truth, alpha, se_type, rng) -> MethodResult:
    # x, y : tableaux NumPy de forme (repeat, n), une ligne par échantillon
    # truth, alpha : nombres ; se_type : texte ("classic", "HC3" ou "HAC")
    # rng : np.random.Generator du flux "method:<id>" (seule source d'aléatoire autorisée)
    ...
\end{lstlisting}

\subsection{Ce qu'une méthode renvoie}
Un objet \texttt{MethodResult} avec quatre blocs, toujours les mêmes :

\begin{table}[H]
  \centering
  \small
  \begin{tabularx}{\textwidth}{@{}L{2.4cm} L{3.6cm} Y@{}}
    \toprule
    \textbf{Bloc} & \textbf{Forme} & \textbf{À quoi il sert} \\
    \midrule
    \texttt{reps} & \code{dict[str, ndarray]}, chaque tableau de forme \texttt{(repeat,)} & Un résultat par échantillon (\texttt{est}, \texttt{low}, \texttt{high}, \texttt{hit}, \texttt{pvalue}, \texttt{diverged}\dots). C'est la matière des agrégateurs et des vues Répétitions et Distribution. \\
    \texttt{example} & \code{dict[str, ndarray]}, tableaux de forme \texttt{(n,)} & Le détail du premier échantillon (\texttt{x}, \texttt{y}, \code{y_hat}, \texttt{resid}, \texttt{leverage}\dots), pour les vues Nuage et Résidus. \\
    \texttt{curves} & \code{dict[str, ndarray]}, facultatif & Des courbes déjà calculées (\texttt{t}, \texttt{theta}, \texttt{paths}\dots), pour les vues Courbes et Répétitions en mode trajectoires. \\
    \texttt{summary} & \code{dict[str, float]}, facultatif & Les seules grandeurs qui ne viennent pas des répétitions (période d'un pendule), lues par l'agrégateur \code{summary(k)}. \\
    \bottomrule
  \end{tabularx}
  \caption{Les quatre blocs renvoyés par une méthode.}
\end{table}

C'est ensuite le moteur (R1) qui calcule les mesures et fabrique les vues à partir de ces blocs. Une méthode ne calcule jamais elle-même une couverture ou un taux de rejet, et ne construit jamais un JSON de vue.

\subsection{Ce que le moteur vérifie}
\begin{itemize}[leftmargin=*, nosep]
  \item \textbf{Au chargement d'un YAML :} chaque chemin lu par une vue (\texttt{bind}) et chaque agrégateur porte sur une clé que la méthode déclare dans \texttt{provides}. Une vue qui lit \code{example.resid} avec une méthode qui ne le fournit pas est refusée avec le numéro de ligne.
  \item \textbf{À l'exécution :} la méthode renvoie exactement les clés déclarées, aux bonnes formes. Une répétition qui échoue numériquement vaut \texttt{NaN} dans \texttt{reps} ; elle est exclue et comptée.
  \item \textbf{Le coût :} \texttt{cost} est une expression en \texttt{repeat}, \texttt{n}, \texttt{p} ou en paramètres du générateur (\code{t_end / dt} pour une intégration), utilisée par le budget de calcul (section~\ref{sec:robustesse}).
  \item \textbf{Les méthodes déterministes} (générateurs \texttt{ode} et \texttt{function}) déclarent \texttt{curves} et \texttt{summary} dans \texttt{provides}, et ne renvoient pas de \texttt{reps} : l'exemple complet est la méthode \texttt{integrate} de l'annexe~\ref{ann:yaml-ode}.
\end{itemize}

\subsection{Les règles d'une méthode}
\begin{itemize}[leftmargin=*, nosep]
  \item \textbf{Fonction pure :} pas d'affichage, pas de fichier, pas d'état global ; l'aléatoire passe uniquement par \texttt{rng}.
  \item \textbf{Vectorisée sur les répétitions :} on traite les 500 échantillons d'un coup (axe 0), sans boucle Python sur les répétitions.
  \item \textbf{Rapide :} moins de 200 ms pour \texttt{repeat = 500} et \texttt{n = 1000}.
  \item \textbf{Testée seule :} \code{tests/test_method_contract.py} appelle chaque méthode enregistrée sur de petites données factices et vérifie les clés déclarées, les formes et les types. R2 peut donc valider ses méthodes le jour, avant que le moteur existe.
\end{itemize}

\newpage
\section{Formats des vues}\label{ann:vues}
Une vue déclare où lire chaque donnée (\texttt{bind}) ; le moteur résout ces chemins et renvoie des séries prêtes à dessiner. C'est le contrat entre le moteur (R1) et \code{PlotlyChart.svelte} (Kilian), et la base des fixtures utilisées dès H+4.

\textbf{Chemins autorisés dans \texttt{bind} :} \code{example.<k>}, \code{reps.<k>}, \code{curves.<k>}, \code{summary.<k>} (sorties d'une méthode), \code{gen.<variable>} (premier échantillon du générateur) et \code{p.<param>}. Un emplacement marqué \og par méthode \fg{} produit une série par méthode ; les autres sont lus sur la méthode principale.

\begin{table}[H]
  \centering
  \footnotesize
  \begin{tabularx}{\textwidth}{@{}L{2.6cm} c L{3.4cm} L{2.9cm} Y@{}}
    \toprule
    \textbf{Type de vue} & \textbf{Prio.} & \textbf{Emplacements obligatoires} & \textbf{Facultatifs} & \textbf{Par méthode} \\
    \midrule
    \texttt{scatter} & P0 & \texttt{x}, \texttt{y} & \texttt{fit}, \texttt{projections}, \texttt{segments} & \texttt{fit} \\
    \texttt{residuals} & P0 & \texttt{x}, \texttt{resid} & \texttt{trend} & \texttt{resid} \\
    \texttt{repetitions} (intervalles) & P0 & \texttt{est}, \texttt{low}, \texttt{high}, \texttt{hit}, \texttt{truth} & aucun & tous sauf \texttt{truth} \\
    \texttt{repetitions} (trajectoires) & P0 & \texttt{paths} & \texttt{truth} & \texttt{paths} \\
    \texttt{distribution} & P0 & \texttt{stat} & \texttt{theory}, \texttt{threshold}, \texttt{observed}, \texttt{h1} & \texttt{stat} \\
    \texttt{curves} & P0 & \texttt{t} & \texttt{series} (liste de \texttt{\{ y, label \}}), \texttt{truth} & \texttt{series} \\
    \texttt{sensitivity} & P0 & \texttt{sweep} (1 axe) & \texttt{measure} & automatique \\
    \code{validity_domain} & P1 & \texttt{sweep} (2 axes) & \texttt{measure} & méthode principale \\
    \code{surface_3d} & P1 & \texttt{x}, \texttt{y}, \texttt{z} & \texttt{fit} & \texttt{fit} \\
    \code{vector_space} & P2 & \texttt{vectors}, \texttt{target} & \texttt{approx} & \texttt{approx} \\
    \code{measures_table} & P0 & aucun & aucun & automatique \\
    \bottomrule
  \end{tabularx}
  \caption{Les emplacements de chaque type de vue. \texttt{measures\_table} est la vue de secours : elle s'affiche quand une vue ne peut pas être construite, ou quand un YAML n'en déclare aucune.}
\end{table}

\textbf{Champs facultatifs d'une vue dans le YAML :} \code{allow_point_injection} (ajout de points, fiche~\ref{fiche:points}), \texttt{inspect} (inspecteur, fiche~\ref{fiche:inspecteur}), \texttt{layers} (liste ordonnée des calques, fiche~\ref{fiche:calques}), \texttt{axes} et le champ libre \texttt{ui}. Dans \texttt{bind}, un emplacement vaut soit un chemin, soit \texttt{\{ path, layer \}} pour le rattacher à un calque. L'emplacement \texttt{segments} du nuage est une liste de \texttt{\{ from, to, layer, label \}} : un segment vertical par observation, de \texttt{from} à \texttt{to}.

\textbf{Champs communs à toute vue renvoyée :} \texttt{id}, \texttt{type}, \texttt{title}, \texttt{axes} (bornes calculées sur la Référence, ou fixées dans le YAML), \texttt{shown} et \texttt{total} (ce qui est affiché et ce qui a été calculé), \texttt{clipped} (points hors cadre), et le champ libre \texttt{ui} transmis tel quel.

\textbf{Champs de chaque série :} \texttt{id} stable (\code{<vue>.<emplacement>.<methode>}), \texttt{method}, \texttt{role}, \texttt{label}, \texttt{layer} facultatif (calques pas à pas), et les données. Grâce aux identifiants stables, \code{Plotly.react()} anime proprement et les calques marchent dans toutes les vues sans code dédié.

\begin{table}[H]
  \centering
  \small
  \begin{tabularx}{\textwidth}{@{}L{3.2cm} L{2.4cm} L{2.6cm} Y L{2.2cm}@{}}
    \toprule
    \textbf{Rôle} & \textbf{Thème sombre} & \textbf{Impression} & \textbf{Trait} & \textbf{Marqueur} \\
    \midrule
    \texttt{method} & \#4ea1ff & \textcolor{ctMethod}{\textbf{\#1D6FD8}} & plein, 2 px & rond \\
    \texttt{correction} & \#3fb950 & \textcolor{ctCorrection}{\textbf{\#1F883D}} & plein, 3 px & losange \\
    \code{alternative_1} & \#d2a8ff & \textcolor{ctAltOne}{\textbf{\#8250DF}} & tirets et points & triangle \\
    \code{alternative_2} & \#56d4dd & \textcolor{ctAltTwo}{\textbf{\#1B7C83}} & points & carré \\
    \texttt{reference} & \#6e7681 & \textcolor{ctBaseline}{\textbf{\#57606A}} & tirets, 1,5 px & rond creux \\
    \texttt{truth} (générateur) & \#f0883e & \textcolor{ctGenerator}{\textbf{\#C95E12}} & plein, 1 px & aucun \\
    \texttt{rupture} & \#f85149 & \textcolor{ctBroken}{\textbf{\#CF222E}} & plein & croix \\
    \texttt{injected} (point posé) & \#e6edf3 & \textbf{\#1F2328} & aucun & étoile \\
    \bottomrule
  \end{tabularx}
  \caption{Les rôles et leurs styles, définis une fois dans \texttt{tokens.css} et \texttt{roles.ts}. Les alternatives reçoivent leur rôle dans l'ordre du YAML.}
\end{table}

\textbf{Limites d'affichage, appliquées par le moteur :} 100 intervalles, 30 trajectoires, 2\,000 points par nuage, 2\,000 points par courbe (sous-échantillonnage régulier). Les mesures sont toujours calculées sur toutes les répétitions.

\newpage
\section{Les mécaniques visuelles en détail}\label{ann:mecaniques}
Une fiche par mécanique de la section~\ref{sec:mecaniques}, toujours avec les mêmes cinq rubriques. Les mécaniques n'ajoutent aucune route : elles s'appuient sur \code{/api/run}, \code{/api/sweep} et sur des champs facultatifs du YAML. Une expérience qui ne déclare pas un champ n'affiche simplement pas la mécanique correspondante.

\subsection{\texorpdfstring{Banc de méthodes \pone}{Banc de méthodes}}\label{fiche:banc}
\begin{fiche}
  \textbf{Ce que voit l'étudiant} & Sous les mesures, une ligne par méthode affichée : nom, valeur $\pm$ marge, badge de verdict, plus une ligne pour la Référence. Un interrupteur par méthode ; le bouton Correction est l'interrupteur de la correction. Dans les vues, chaque méthode affichée a sa propre série, avec le style de son rôle (annexe~\ref{ann:vues}). \\
  \midrule
  \textbf{Activé dans le YAML par} & La liste \texttt{methods} : 1 méthode principale, au plus 1 correction, au plus 2 alternatives. Le \emph{coût} d'une méthode robuste se lit sur une mesure secondaire sans cible que l'expérience déclare : largeur moyenne des IC (\texttt{width}), puissance d'un test, variance d'un estimateur. \\
  \midrule
  \textbf{Données du moteur} & Rien de spécifique : toutes les méthodes sont calculées à chaque exécution, et chaque mesure renvoie \texttt{values}, \code{mc_se}, \texttt{states}, \texttt{texts} et \texttt{baseline} par méthode (annexe~\ref{ann:api}). \\
  \midrule
  \textbf{Côté interface} & \code{MethodBench.svelte}. Les méthodes affichées sont dans l'état d'affichage (\code{methods_on}, \texttt{on=} dans l'URL). Afficher ou masquer une méthode ne fait aucun appel au serveur. \\
  \midrule
  \textbf{Règles et cas limites} & Le bouton Correction seul est \pzero{} ; le banc complet (alternatives, ligne de coût) est \pone{}. Par défaut, seule la méthode principale est affichée ; la démo s'ouvre ainsi, puis on clique sur Correction. Les alternatives prennent les rôles \code{alternative_1} puis \code{alternative_2} dans l'ordre du YAML. Sans correction déclarée, le bouton Correction n'apparaît pas. À $h = 0$, la ligne de la mesure de coût compare chaque méthode à la méthode principale (\og puissance 63\,\% contre 67\,\% \fg{} pour Wilcoxon face au test de Student sous loi normale, à $n = 25$). \\
\end{fiche}

\subsection{\texorpdfstring{Contre-exemples en un clic \pzero}{Contre-exemples en un clic}}\label{fiche:contre}
\begin{fiche}
  \textbf{Ce que voit l'étudiant} & Une barre \og Cas limites \fg{} au-dessus des contrôles, avec un bouton par contre-exemple. Un clic règle tous les contrôles d'un coup et affiche la phrase \texttt{explain} dans un bandeau. \\
  \midrule
  \textbf{Activé dans le YAML par} & \texttt{counterexamples} : une liste de \texttt{\{ id, label, explain, state \}}. Pour Anscombe, le paramètre \code{match_summary} active le calage, et l'hypothèse \texttt{relation} ou \code{add_leverage} change la forme du nuage (annexe~\ref{ann:yaml}). \\
  \midrule
  \textbf{Données du moteur} & Aucun calcul dédié : un contre-exemple est un état comme un autre. La fonction \code{ols_match(x, y, a, b, r2)} de la liste blanche recale chaque répétition pour obtenir exactement l'ordonnée à l'origine $a$, la pente $b$ et le $R^2$ demandés. Le test 4 de la section~\ref{sec:tests} vérifie que chaque contre-exemple rompt au moins une mesure. \\
  \midrule
  \textbf{Côté interface} & \code{CounterexampleBar.svelte}, qui écrit l'état de calcul complet dans \code{state.ts}. \\
  \midrule
  \textbf{Règles et cas limites} & Les valeurs absentes de \texttt{state} reviennent à l'état sain ou à la valeur par défaut ; la graine est conservée ; les points ajoutés sont effacés. L'URL change, donc un contre-exemple se partage comme n'importe quel état. La barre est masquée dans l'Entraînement. \\
\end{fiche}

\subsection{\texorpdfstring{Balayage automatique \pone}{Balayage automatique}}\label{fiche:balayage}
\begin{fiche}
  \textbf{Ce que voit l'étudiant} & Un bouton Lecture à côté de chaque contrôle animable. Pour une hypothèse : elle monte de 0 à 1, les vues bougent, l'animation marque une pause de 2 secondes au seuil critique (le verdict passe à Fragile), puis la Correction s'active à la fin. Pour $n$ : il parcourt ses valeurs de balayage, et une phrase conclut \og Le grand échantillon répare \fg{} ou \og Le grand échantillon ne répare pas \fg{}, selon le verdict à la première et à la dernière valeur. \\
  \midrule
  \textbf{Activé dans le YAML par} & \code{auto_sweep} : la liste des contrôles animables, dans l'ordre des boutons, par exemple \texttt{[h.hetero, p.n]}. Un paramètre doit déclarer \code{sweep_values}. \\
  \midrule
  \textbf{Données du moteur} & Un appel à \code{/api/sweep} pour connaître \code{critical.warning_at}, puis un appel à \code{/api/run} par valeur (11 pour une hypothèse continue), mis en cache côté serveur après le premier passage. \\
  \midrule
  \textbf{Côté interface} & \code{AutoSweepPlayer.svelte}. Toutes les images sont préchargées avant le départ (barre de progression), puis jouées à 600 ms par image. \\
  \midrule
  \textbf{Règles et cas limites} & Le moindre clic ou glissement de l'étudiant arrête l'animation et garde l'état affiché. La pause tombe sur la première valeur de la grille supérieure ou égale à \code{warning_at} ; sans seuil critique, il n'y a pas de pause. Pendant une animation sur $n$, les axes de la première image sont conservés. La Correction n'est activée à la fin que pour une hypothèse, jamais pour $n$. \\
\end{fiche}

\subsection{\texorpdfstring{Ajout de points par clic \pone}{Ajout de points par clic}}\label{fiche:points}
\begin{fiche}
  \textbf{Ce que voit l'étudiant} & Un clic dans une zone vide du graphique pose un point (marqueur étoile) et toutes les vues se recalculent. Un clic sur un point posé le retire. Un compteur indique \og 3 / 20 points \fg{}. \\
  \midrule
  \textbf{Activé dans le YAML par} & Dans la vue, \code{allow_point_injection: true}. Dans le générateur, \texttt{inject} associe les axes de la vue aux variables : \texttt{\{ x: x, y: y \}} pour un nuage, \texttt{\{ x: obs \}} pour une vue à une dimension (pour comparer moyenne et médiane, le clic pose une valeur sur l'axe horizontal d'un histogramme). Le chargeur refuse une vue qui autorise l'ajout sans \texttt{inject} dans le générateur. \\
  \midrule
  \textbf{Données du moteur} & Les points arrivent dans \code{extra_points} (état de calcul). Ils sont ajoutés à toutes les répétitions après le calcul des variables, jamais à la Référence. Ils font partie de \texttt{example}, donc du nuage, des résidus et de l'inspecteur. \\
  \midrule
  \textbf{Côté interface} & Événement \code{plotly_click} dans \code{PlotlyChart.svelte}, converti en coordonnées de données ; les points sont écrits dans l'URL (\texttt{pts=}). Rôle \texttt{injected} pour le style (annexe~\ref{ann:vues}). \\
  \midrule
  \textbf{Règles et cas limites} & 20 points au plus, ensuite le clic affiche un message. Désactivé dans l'Entraînement, réactivé après \og Ouvrir dans le Laboratoire \fg{}. Un contre-exemple efface les points. Partage du clic avec l'inspecteur : clic sur un point de l'échantillon = inspecteur, clic sur un point posé = retrait, clic dans le vide = ajout. Un vecteur cible en algèbre linéaire n'utilise pas cette mécanique : c'est un paramètre (expérience \ptwo{}). \\
\end{fiche}

\subsection{\texorpdfstring{Journal des tirages \pone}{Journal des tirages}}\label{fiche:journal}
\begin{fiche}
  \textbf{Ce que voit l'étudiant} & À côté de \og Nouveau tirage \fg{}, une rangée de pastilles : verte si le test ne rejette pas, rouge s'il rejette, avec la $p$-value au survol. Une pastille est le test sur le \emph{premier échantillon} du tirage (la répétition 0), c'est-à-dire ce que ferait un étudiant avec un seul jeu de données. Une phrase compare au taux calculé sur toutes les répétitions : \og 15 tirages sous $H_0$ : 1 faux positif (6,7\,\%). Sur les 500 tirages simulés : 5,2\,\%. Théorie : 5\,\%. \fg{} \\
  \midrule
  \textbf{Activé dans le YAML par} & \code{draw_log}, au niveau de l'expérience : \texttt{\{ reject: reps.reject, pvalue: reps.pvalue, reject\_label: "faux positif" \}}. Le libellé dit ce qu'est un rejet dans cette expérience (\og faux positif \fg{} sous $H_0$, \og détection \fg{} sous $H_1$). La méthode principale doit fournir \texttt{reps.reject} et \texttt{reps.pvalue} ; le chargeur le vérifie. \\
  \midrule
  \textbf{Données du moteur} & \code{draw_log} dans \texttt{RunResult} : \texttt{\{ reject, pvalue \}} lus sur la répétition 0 de la méthode principale ; \texttt{null} si l'expérience ne déclare pas le champ. \\
  \midrule
  \textbf{Côté interface} & \code{DrawLog.svelte}. Le journal vit dans la mémoire de la page (ni URL, ni \texttt{localStorage}). \\
  \midrule
  \textbf{Règles et cas limites} & Le journal se vide dès qu'un contrôle autre que la graine change, pour ne jamais mélanger des tirages faits dans des conditions différentes. 30 pastilles visibles ; les compteurs portent sur tout le journal. \\
\end{fiche}

\subsection{\texorpdfstring{Inspecteur d'observation \ptwo}{Inspecteur d'observation}}\label{fiche:inspecteur}
\begin{fiche}
  \textbf{Ce que voit l'étudiant} & Un clic sur un point du nuage ou des résidus l'entoure et ouvre sa fiche : $(x_i, y_i)$, prédiction $\hat{y}_i$, résidu $\hat{\varepsilon}_i$, levier $h_{ii}$, distance de Cook $D_i$, et part du point dans la somme des carrés des résidus, $\hat{\varepsilon}_i^2 / \sum_j \hat{\varepsilon}_j^2$. Les seuils usuels sont signalés : $h_{ii} > 2p/n$ et $D_i > 4/n$. \\
  \midrule
  \textbf{Activé dans le YAML par} & Dans la vue, \texttt{inspect: true}. \\
  \midrule
  \textbf{Données du moteur} & La méthode principale fournit \code{example.leverage} et \code{example.cooks_d} ; le chargeur le vérifie quand une vue déclare \texttt{inspect}. Le reste (prédiction, résidu, part dans la somme) est lu ou calculé par l'interface à partir de \texttt{example}. \\
  \midrule
  \textbf{Côté interface} & \code{Inspector.svelte}, ouvert par \code{plotly_click} sur un point existant de l'échantillon. \\
  \midrule
  \textbf{Règles et cas limites} & La fiche se met à jour quand l'état change et se ferme si le point disparaît. Un point posé à la main s'inspecte comme les autres. Partage du clic avec l'ajout de points : voir la fiche~\ref{fiche:points}. \\
\end{fiche}

\subsection{\texorpdfstring{Calques pas à pas \ptwo}{Calques pas à pas}}\label{fiche:calques}
\begin{fiche}
  \textbf{Ce que voit l'étudiant} & Un curseur d'étapes au-dessus de la vue ; chaque étape ajoute une couche à la figure. Régression et ANOVA : le nuage, puis les écarts à la moyenne $y_i - \bar{y}$ ($SCT$), puis la droite et les écarts expliqués $\hat{y}_i - \bar{y}$ ($SCE$), puis les résidus $y_i - \hat{y}_i$ ($SCR$), avec les trois sommes et $SCT = SCE + SCR$ vérifié en direct. Tests : la loi sous $H_0$, puis les zones de rejet, puis la statistique observée et la $p$-value, puis la loi sous $H_1$ avec $\beta$ et la puissance. \\
  \midrule
  \textbf{Activé dans le YAML par} & La vue déclare ses calques dans l'ordre (\texttt{layers}, une liste de \texttt{\{ id, label \}}), et chaque emplacement de \texttt{bind} peut être rattaché à un calque : \texttt{\{ path, layer \}}. Exemple complet dans l'annexe~\ref{ann:yaml} (vue \texttt{nuage}). \\
  \midrule
  \textbf{Données du moteur} & Le moteur recopie l'identifiant du calque sur chaque série. Pour la régression, la méthode fournit \code{example.y_mean} (la moyenne $\bar{y}$ répétée $n$ fois), utilisée par l'emplacement \texttt{segments}. \\
  \midrule
  \textbf{Côté interface} & \code{LayerStepper.svelte}. L'étape $k$ affiche les calques 1 à $k$ ; les séries sans calque restent toujours visibles. Aucun appel au serveur. Les sommes $SCT$, $SCE$ et $SCR$ sont calculées par l'interface à partir des segments affichés. \\
  \midrule
  \textbf{Règles et cas limites} & Une vue sans \texttt{layers} s'affiche en entier, sans curseur d'étapes. Un identifiant de calque utilisé dans \texttt{bind} mais absent de \texttt{layers} est refusé au chargement. \\
\end{fiche}

\newpage
\section{Arborescence du dépôt}\label{ann:arbo}
Chaque bloc a un propriétaire : R1 (moteur), R2 (maths et contenu), R3 (Kilian, interface), R4 (écrans secondaires et qualité). L'assistant IA de chacun ne modifie que ses propres fichiers. Les fichiers marqués \textbf{[Contrat]} décrivent les échanges entre les parties : ils ne changent pas pendant le hackathon.

\renewcommand{\arraystretch}{1.15}
\begin{longtable}{@{}L{6.4cm} >{\raggedright\arraybackslash}p{\dimexpr\textwidth-7.9cm-4\tabcolsep\relax} L{1.5cm}@{}}
  \toprule
  \textbf{Dossier / Fichier} & \textbf{Rôle} & \textbf{Resp.} \\
  \midrule
  \endfirsthead
  \toprule
  \textbf{Dossier / Fichier} & \textbf{Rôle} & \textbf{Resp.} \\
  \midrule
  \endhead
  \bottomrule
  \endfoot

  \rowcolor{ctRowShade}
  \multicolumn{3}{@{}l@{}}{\textbf{Racine du dépôt}} \\
  \texttt{CLAUDE.md} & Règles des assistants IA (Svelte 5, vouvoiement, registres, dossiers autorisés). & R3 \\
  \texttt{pyproject.toml}, \texttt{uv.lock} & Dépendances Python (\texttt{uv}) aux versions fixées, configuration \texttt{pytest} et \texttt{ruff}. & R1 \\
  \texttt{Dockerfile} & Image du moteur, identique en local et en ligne (\code{OMP_NUM_THREADS=1}). & R1 \\
  \code{supabase/schema.sql} & Table \code{shared_experiments} (liens de partage). & R3 \\
  \code{scripts/gen_types.sh} & Produit \code{types.gen.ts} à partir de \code{contracts.py}. & R3 \\
  \code{scripts/make_fixtures.py} & Régénère les fixtures de l'interface à partir de vrais résultats. & R1 \\

  \rowcolor{ctRowShade}
  \multicolumn{3}{@{}l@{}}{\textbf{Serveur FastAPI (\code{backend/})}} \\
  \code{backend/main.py} & Point d'entrée, sert aussi le build de l'interface. & R1 \\
  \code{backend/contracts.py} & \textbf{[Contrat]} Modèles Pydantic de l'API (\texttt{RunRequest}, \texttt{RunResult}, \texttt{SweepRequest}, \texttt{SweepResult}\dots). & R1 \\
  \code{backend/api/routes_run.py} & \code{/api/run}, \code{/api/sweep}. & R1 \\
  \code{backend/api/routes_experiments.py} & \code{/api/experiments}, \code{/api/validate}. & R1 \\
  \code{backend/api/routes_training.py} & \code{/api/training/next}, \code{/api/training/session}. & R1 \\
  \code{backend/engine/training.py} & Fabrication et vérification des séries, réserve par filtre, remplacement en arrière-plan quand le serveur est libre. & R1 \\
  \code{backend/engine/schema.py} & \textbf{[Contrat]} Schéma YAML v1 (\texttt{ExperimentSpec}) et ses vérifications au chargement. & R1 \\
  \code{backend/engine/method_contract.py} & \textbf{[Contrat]} Décorateur, \texttt{provides}, \texttt{cost} et \texttt{MethodResult} (annexe~\ref{ann:methodes}). & R2 \\
  \code{backend/engine/expr.py} & Évaluateur sécurisé (\texttt{asteval}), liste blanche et espace de noms. & R1 \\
  \code{backend/engine/rng.py} & Un flux aléatoire par variable, dérivé de la graine et du nom. & R1 \\
  \code{backend/engine/laws.py} & Registre des lois, tirées depuis des uniformes communes. & R2 \\
  \code{backend/engine/aggregators.py} & Agrégateurs et erreurs Monte Carlo. & R2 \\
  \code{backend/engine/sanitize.py} & \texttt{NaN}/\texttt{Inf} vers \texttt{null}, avertissements en français à code stable. & R1 \\
  \code{backend/engine/runner.py} & \texttt{run()}, Référence, balayage, verdicts, cache et budget. & R1 \\
  \code{backend/engine/loader.py} & Lecture du YAML et erreurs par numéro de ligne. & R1 \\
  \code{backend/engine/generators/} & \texttt{random.py}, \texttt{ode.py} (P0) ; \texttt{function.py} (P2). & R1 \\
  \code{backend/methods/} & \texttt{stats.py}, \code{tests_stat.py}, \texttt{probability.py}, \texttt{dynamics.py}, \texttt{analysis.py}. & R2 \\

  \rowcolor{ctRowShade}
  \multicolumn{3}{@{}l@{}}{\textbf{Expériences (\code{experiments/})}} \\
  \code{experiments/regression_lineaire.yaml} & \textbf{[P0]} Fichier de référence (écrit ensemble). & R2 \\
  \code{experiments/test_student.yaml} & \textbf{[P0]} Lois symétriques, risque $\alpha$ et puissance, Wilcoxon. & R2 \\
  \code{experiments/pendule_euler.yaml} & \textbf{[P0]} EDO, pas de temps, Euler et RK4 (modèle du générateur \texttt{ode}). & R2 \\
  \code{experiments/lgn_tcl.yaml} & \textbf{[P1]} Trajectoires, loi asymétrique, Cauchy, moyenne et médiane. & R2 \\
  \code{experiments/regression_3d.yaml} & \textbf{[P2]} Plan 3D, multicolinéarité, Ridge. & R2 \\
  \code{experiments/analyse_tfa.yaml} & \textbf{[P2]} $f(x)$ libre, continuité, intégrale. & R2 \\
  \code{experiments/pendule_petits_angles.yaml} & \textbf{[P2]} Approximation $\sin\theta \approx \theta$, formule de Borda. & R2 \\

  \rowcolor{ctRowShade}
  \multicolumn{3}{@{}l@{}}{\textbf{Interface Vite + Svelte 5 (\code{frontend/})}} \\
  \code{frontend/src/App.svelte} & Navigation (Accueil, Laboratoire, Entraînement, Bilan, Import). & R3 \\
  \code{frontend/src/styles/tokens.css} & Thème sombre et palette de sens. & R3 \\
  \code{frontend/src/lib/types.gen.ts} & Types TypeScript générés depuis \code{contracts.py} ; jamais modifié à la main. & R3 \\
  \code{frontend/src/lib/roles.ts} & Couleur, trait et marqueur de chaque rôle. & R3 \\
  \code{frontend/src/lib/equation.ts} & Règle de couleur des termes de l'équation. & R3 \\
  \code{frontend/src/lib/state.ts} & État de calcul et état d'affichage, forme canonique, lecture et écriture de l'URL. & R3 \\
  \code{frontend/src/lib/executor.ts} & Client API unique (regroupement 100 ms, annulation), adresse lue dans \code{VITE_API_URL}. & R3 \\
  \code{frontend/src/lib/share.ts} & Liens de partage (URL et Supabase). & R3 \\
  \code{frontend/src/lib/calibration.ts} & Crans, score de Brier, courbe, profil démo. & R4 \\
  \code{frontend/src/fixtures/} & Résultats JSON pour développer sans moteur (écrits à la main jusqu'à H+4, puis régénérés). & R1 \\
  \code{frontend/src/components/_Modele.svelte} & Composant modèle pour les écrans de la nuit. & R3 \\
  \code{frontend/src/components/PlotlyChart.svelte} & Composant graphique unique (rôles, axes stables, calques). & R3 \\
  \code{frontend/src/components/DynamicEquation.svelte} & Équation KaTeX à termes nommés. & R3 \\
  \code{frontend/src/components/lab/} & \code{LabView}, \code{ControlPanel}, \code{UnderstandPanel}, \code{ChallengesTab}, et les mécaniques visuelles : \code{MethodBench}, \code{CounterexampleBar}, \code{AutoSweepPlayer}, \code{DrawLog}, \code{Inspector}, \code{LayerStepper} (annexe~\ref{ann:mecaniques}). & R3 \\
  \code{frontend/src/components/training/} & \code{TrainingView}, \code{ValidationPanel}, \code{ConfidencePicker}. & R3 \\
  \code{frontend/src/components/home/} & \code{Home}, \code{Gallery}. & R4 \\
  \code{frontend/src/components/summary/} & \code{CalibrationSummary} (écran Bilan). & R4 \\
  \code{frontend/src/components/ImportView.svelte} & Prompt à copier, éditeur, validation, lien de partage. & R4 \\

  \rowcolor{ctRowShade}
  \multicolumn{3}{@{}l@{}}{\textbf{Tests (\code{tests/})}} \\
  \code{tests/test_engine.py} & \texttt{expr}, \texttt{rng}, \texttt{sanitize}, \texttt{runner}, générateurs. & R1 \\
  \code{tests/test_loader.py} & Le YAML de référence est accepté, les YAML volontairement faux sont refusés à la bonne ligne. & R4 \\
  \code{tests/test_experiments.py} & Les dix vérifications scientifiques de chaque YAML. & R2 \\
  \code{tests/test_properties.py} & États tirés au hasard : aucune exception, aucun \texttt{NaN}. & R1 \\
  \code{tests/test_api.py} & Tests d'intégration des routes, et concordance entre \code{types.gen.ts} et \code{contracts.py}. & R1 \\
  \code{tests/test_method_contract.py} & Conformité de chaque méthode à ce qu'elle déclare. & R2 \\
\end{longtable}
\renewcommand{\arraystretch}{1.0}

\newpage
\section{Le score de Brier en détail}\label{ann:brier}

\subsection{Définition}
Pour $N$ réponses, avec $p_i$ la probabilité associée au cran choisi et $y_i \in \{0,1\}$ le résultat :
\[
  \mathrm{Brier} = \frac{1}{N} \sum_{i=1}^{N} (p_i - y_i)^2 \in [0, 1].
\]
Avec $k$ choix, les crans valent $p_{\text{devine}} = 1/k$, $p_{\text{sûr}} = \tfrac{1}{2}\left(1/k + 0{,}95\right)$ et $p_{\text{certain}} = 0{,}95$.

\subsection{Pénalité par réponse}
\begin{table}[H]
  \centering
  \small
  \begin{tabular}{@{}llccc@{}}
    \toprule
    \textbf{Bouton} & \textbf{Question} & \textbf{$p$} & \textbf{Réponse juste} & \textbf{Réponse fausse} \\
    \midrule
    Certain & toutes & 95\,\% & 0,0025 & 0,9025 \\
    \midrule
    Plutôt sûr & binaire ($k=2$) & 72,5\,\% & 0,0756 & 0,5256 \\
    Plutôt sûr & 3 choix ($k=3$) & 64,2\,\% & 0,1284 & 0,4117 \\
    Plutôt sûr & 4 choix ($k=4$) & 60\,\% & 0,16 & 0,36 \\
    \midrule
    Je devine & binaire ($k=2$) & 50\,\% & 0,25 & 0,25 \\
    Je devine & 3 choix ($k=3$) & 33,3\,\% & 0,4444 & 0,1111 \\
    Je devine & 4 choix ($k=4$) & 25\,\% & 0,5625 & 0,0625 \\
    \bottomrule
  \end{tabular}
  \caption{Contribution d'une réponse au score de Brier, pour chaque cran.}
\end{table}

\subsection{Trois profils types}
Les trois profils portent sur des QCM à 4 choix.
\begin{itemize}[leftmargin=*, nosep]
  \item \textbf{Lucide :} 60\,\% de réussite, clique \og Plutôt sûr \fg{} (60\,\%) : $0{,}6 \times 0{,}16 + 0{,}4 \times 0{,}36 = 0{,}24$.
  \item \textbf{Trop sûr de lui :} 60\,\% de réussite, clique \og Certain \fg{} (95\,\%) : $0{,}6 \times 0{,}0025 + 0{,}4 \times 0{,}9025 \approx 0{,}36$.
  \item \textbf{Ignorant mais honnête :} répond au hasard (25\,\% de réussite) et clique \og Je devine \fg{} (25\,\%) : $0{,}25 \times 0{,}5625 + 0{,}75 \times 0{,}0625 \approx 0{,}19$.
\end{itemize}
Le score récompense donc l'honnêteté sur ce que l'on sait, et pas seulement le nombre de bonnes réponses : l'étudiant qui ne sait rien mais l'assume obtient un meilleur score que celui qui en sait plus mais se surestime.

\subsection{La courbe de calibration}
On regroupe les réponses par cran de confiance. Pour chaque cran $c$, on calcule le taux de réussite observé $\bar{o}_c$ et on le place face à la probabilité moyenne annoncée dans ce cran, $p_c$. Un étudiant bien calibré a ses points sur la diagonale ; un point sous la diagonale signale un excès de confiance. Regrouper par \emph{bouton} (et non par pourcentage) permet de mélanger questions binaires et à 4 choix sur un même graphique. La courbe n'est affichée qu'à partir d'une vingtaine de réponses, pour éviter un graphique trop bruité.

\subsection{Pour aller plus loin : la décomposition de Murphy}
On regroupe ici les réponses par \emph{valeur} de $p_i$ (elles n'en prennent qu'un petit nombre : 0,25, 0,33, 0,5, 0,6, 0,64, 0,725, 0,95). Avec $n_c$ réponses dans le groupe $c$, de probabilité commune $p_c$, et $\bar{o}$ le taux de réussite global, le score se décompose exactement en trois termes :
\[
  \mathrm{Brier} =
  \underbrace{\frac{1}{N}\sum_{c} n_c\,(p_c - \bar{o}_c)^2}_{\text{calibration (à minimiser)}}
  \; - \;
  \underbrace{\frac{1}{N}\sum_{c} n_c\,(\bar{o}_c - \bar{o})^2}_{\text{résolution (à maximiser)}}
  \; + \;
  \underbrace{\bar{o}\,(1 - \bar{o})}_{\text{incertitude}} .
\]
La \emph{calibration} mesure si la confiance colle à la réussite ; la \emph{résolution} mesure si l'étudiant sait distinguer les questions qu'il maîtrise de celles qu'il ne maîtrise pas ; l'\emph{incertitude} dépend seulement de la difficulté des questions. C'est une bonne réponse si le jury demande pourquoi ce score plutôt que le simple pourcentage de bonnes réponses.

\end{appendices}

\end{document}
